%%%%%%%%%%%%%%%%%%%%%%%%%%%%%%%%%%%%%%%%%%%%%%%%%%%%%%%%%%%%%%%% 
% 
% Apuntes de la asignatura Geometría III.
% Doble Grado de Informática y Matemáticas.
% Universidad de Granada.
% Curso 2016/17.
% 
% 
% Colaboradores:
% Javier Sáez (@fjsaezm)
% Daniel Pozo (@danipozodg)
% Pedro Bonilla (@pedrobn23)
% Guillermo Galindo(@guillegalor)
% Antonio Coín (@antcc)
% Sofía Almeida (@SofiaAlmeida)
% Miguel Lentisco (@AlfaOmegaX)
% 
% Agradecimientos:
% Andrés Herrera (@andreshp) y Mario Román (@M42) por
% las plantillas base.
% 
% Sitio original:
% https://github.com/libreim/apuntesDGIIM/
% 
% Licencia:
% CC BY-NC-SA 4.0 (https://creativecommons.org/licenses/by-nc-sa/4.0/)
% 
%%%%%%%%%%%%%%%%%%%%%%%%%%%%%%%%%%%%%%%%%%%%%%%%%%%%%%%%%%%%%%% 


% ------------------------------------------------------------------------------
% ACKNOWLEDGMENTS
% ------------------------------------------------------------------------------

%%%%%%%%%%%%%%%%%%%%%%%%%%%%%%%%%%%%%%%%%%%%%%%%%%%%%%%%%%%%%%%%%%%%%%%% 
% Plantilla básica de Latex en Español.
% 
% Autor: Andrés Herrera Poyatos (https://github.com/andreshp)
% 
% Es una plantilla básica para redactar documentos. Utiliza el paquete  fancyhdr para darle un
% estilo moderno pero serio.
% 
% La plantilla se encuentra adaptada al español.
% 
%%%%%%%%%%%%%%%%%%%%%%%%%%%%%%%%%%%%%%%%%%%%%%%%%%%%%%%%%%%%%%%%%%%%%%%%% 

%%% 
% Plantilla de Trabajo
% Modificación de una plantilla de Latex de Frits Wenneker para adaptarla
% al castellano y a las necesidades de escribir informática y matemáticas.
% 
% Editada por: Mario Román
% 
% License:
% CC BY-NC-SA 3.0 (http://creativecommons.org/licenses/by-nc-sa/3.0/)
%%% 

%%%%%%%%%%%%%%%%%%%%%%%%%%%%%%%%%%%%%%%% 
% Short Sectioned Assignment
% LaTeX Template
% Version 1.0 (5/5/12)
% 
% This template has been downloaded from:
% http://www.LaTeXTemplates.com
% 
% Original author:
% Frits Wenneker (http://www.howtotex.com)
% 
% License:
% CC BY-NC-SA 3.0 (http://creativecommons.org/licenses/by-nc-sa/3.0/)
% 
%%%%%%%%%%%%%%%%%%%%%%%%%%%%%%%%%%%%%%%%% 


% Tipo de documento y opciones.
%\documentclass[11pt, a4paper, twoside]{article} % Usar para imprimir
\documentclass[11pt, a4paper]{article}


% ---------------------------------------------------------------------------
% PAQUETES
% ---------------------------------------------------------------------------

% Idioma y codificación para Español.
\usepackage[utf8]{inputenc}
\usepackage[spanish, es-tabla, es-lcroman, es-noquoting]{babel}
\selectlanguage{spanish}
% \usepackage[T1]{fontenc}

% Fuente utilizada.
\usepackage{courier}    % Fuente Courier.
\usepackage{microtype}  % Mejora la letra final de cara al lector.

% Diseño de página.
\usepackage{fancyhdr}   % Utilizado para hacer títulos propios.
\usepackage{titlesec} 	% Utilizado para hacer títulos propios.
\usepackage{lastpage}   % Referencia a la última página.
\usepackage{extramarks} % Marcas extras. Utilizado en pie de página y cabecera.
\usepackage[parfill]{parskip}    % Crea una nueva línea entre párrafos.
\usepackage{geometry}            % Geometría de las páginas.

% Símbolos y matemáticas.
\usepackage{amssymb, amsmath, amsthm, amsfonts, amscd}
\usepackage{upgreek}

\usepackage{mdframed}

% Otros.
\usepackage{enumitem}   % Listas mejoradas.
\usepackage[hidelinks]{hyperref}
\usepackage{pgfplots}
\usepackage{graphicx}   % Gráficos.

% Dibujos Tikz
\usepackage{pgf,tikz}
\usepackage{mathrsfs}
\usetikzlibrary{arrows,calc,intersections,through,backgrounds}

% Fuentes personalizadas
\usepackage[scaled=.85]{FiraSans}
\usepackage{ebgaramond-maths}
\usepackage[T1]{fontenc}

% Código para ajustar las fuentes matemáticas al estilo del texto
% que le rodea

\DeclareMathVersion{sans}
\SetSymbolFont{operators}{sans}{OT1}{cmbr}{m}{n}
\SetSymbolFont{letters}{sans}{OML}{cmbr}{m}{it}
\SetSymbolFont{symbols}{sans}{OMS}{cmbrs}{m}{n}
\SetMathAlphabet{\mathit}{sans}{OT1}{cmbr}{m}{sl}
\SetMathAlphabet{\mathbf}{sans}{OT1}{cmbr}{bx}{n}
\SetMathAlphabet{\mathtt}{sans}{OT1}{cmtl}{m}{n}
\SetSymbolFont{largesymbols}{sans}{OMX}{iwona}{m}{n}

\DeclareMathVersion{boldsans}
\SetSymbolFont{operators}{boldsans}{OT1}{cmbr}{b}{n}
\SetSymbolFont{letters}{boldsans}{OML}{cmbrm}{b}{it}
\SetSymbolFont{symbols}{boldsans}{OMS}{cmbrs}{b}{n}
\SetMathAlphabet{\mathit}{boldsans}{OT1}{cmbr}{b}{sl}
\SetMathAlphabet{\mathbf}{boldsans}{OT1}{cmbr}{bx}{n}
\SetMathAlphabet{\mathtt}{boldsans}{OT1}{cmtl}{b}{n}
\SetSymbolFont{largesymbols}{boldsans}{OMX}{iwona}{bx}{n}

\newif\IfInSansMode
\let\oldsf\sffamily
\renewcommand*{\sffamily}{\oldsf\mathversion{sans}\InSansModetrue}
\let\oldmd\mdseries
\renewcommand*{\mdseries}{\oldmd\IfInSansMode\mathversion{sans}\fi\relax}
\let\oldbf\bfseries
\renewcommand*{\bfseries}{\oldbf\IfInSansMode\mathversion{boldsans}\else%
   \mathversion{bold}\fi\relax}
\let\oldnorm\normalfont
\renewcommand*{\normalfont}{\oldnorm\InSansModefalse\mathversion{normal}}
\let\oldrm\rmfamily
\renewcommand*{\rmfamily}{\oldrm\InSansModefalse\mathversion{normal}}

% Colores

\definecolor{50}{HTML}{FFF3E0}
\definecolor{300}{HTML}{FFB74D}
\definecolor{500}{HTML}{FF9800}
\definecolor{700}{HTML}{F57C00}
\definecolor{900}{HTML}{E65100}

% ---------------------------------------------------------------------------
% OPCIONES PERSONALIZADAS
% ---------------------------------------------------------------------------

% Redefinir letra griega épsilon.
\let\epsilon\upvarepsilon


% Formato de texto.
\linespread{1.3}            % Espaciado entre líneas.
\setlength\parindent{0pt}   % No indentar el texto por defecto.
\setlist{leftmargin=.5in}   % Indentación para las listas.

% Estilo de página.
\pagestyle{fancy}
\fancyhf{}
\geometry{left=3cm,right=3cm,top=3cm,bottom=3cm}   % Márgenes y cabecera.

% Estilo de las cabeceras

\titleformat{\section}
  {\Large\bfseries\sffamily}{\thesection}{1em}{}
\titleformat{\subsection}
  {\large\sffamily}{\thesubsection}{1em}{}
\titleformat{\subsubsection}
  {\normalsize\sffamily}{\thesubsubsection}{1em}{}

% Estilo de enlaces
\hypersetup{
  % hidelinks = true,   % Oculta todos los enlaces.
  colorlinks = true,   % Muestra todos los enlaces, sin bordes alrededor.
  linkcolor=black,     % Color de enlaces genéricos
  citecolor={blue!50!black},   % Color de enlaces de referencias
  urlcolor={blue!80!black}     % Color de enlaces de URL
}

% Ruta donde buscar gráficos
\graphicspath{{../Recursos/Plantillas/}, {Recursos/Plantillas/}}

% Redefinir entorno de demostración (reducir espacio superior)
\makeatletter
\renewenvironment{proof}[1][\proofname] {\vspace{-15pt}\par\pushQED{\qed}\normalfont\topsep6\p@\@plus6\p@\relax\trivlist\item[\hskip\labelsep\it#1\@addpunct{.}]\ignorespaces}{\popQED\endtrivlist\@endpefalse}
\makeatother

% Aumentar el tamaño del interlineado
\linespread{1.3}


% ---------------------------------------------------------------------------
% COMANDOS PERSONALIZADOS
% ---------------------------------------------------------------------------

% Valor absoluto: \abs{}
\providecommand{\abs}[1]{\lvert#1\rvert}

% Fracción grande: \ddfrac{}{}
\newcommand\ddfrac[2]{\frac{\displaystyle #1}{\displaystyle #2}}

% Texto en negrita en modo matemática: \bm{}
\newcommand{\bm}[1]{\boldsymbol{#1}}

% Línea horizontal.
\newcommand{\horrule}[1]{\rule{\linewidth}{#1}}

% Renombrar la R de los reales

\newcommand{\R}{\mathbb{R}}

% Flechitas de los vectores de tamaño apropiado
\renewcommand{\vec}{\overrightarrow}

% ---------------------------------------------------------------------------
% CABECERA Y PIE DE PÁGINA
% ---------------------------------------------------------------------------

\renewcommand{\sectionmark}[1]{%
\markboth{#1}{}}

\renewcommand{\subsectionmark}[1]{%
\markright{#1}{}}

%\addtolength{\headheight}{4ex}
\renewcommand{\headrulewidth}{0pt}

\addtolength{\headwidth}{\marginparsep}
%\addtolength{\headwidth}{\marginparwidth}

\fancypagestyle{section}{%
  \fancyhead{}
  %\addtolength{\headheight}{-10ex}
  %\renewcommand{\headheight}{0pt}%
  %\setlength{\footskip}{-48pt}%
  \fancyfoot[LE,RO]{\Large\sffamily\thepage}
  \renewcommand{\headrulewidth}{0pt}%
  \renewcommand{\footrulewidth}{0pt}%
}

\let\originalsection\section
\RenewDocumentCommand{\section}{som}{%
  \IfBooleanTF{#1}
    {\originalsection*{#3}}
    {\IfNoValueTF{#2}
      {\originalsection{#3}}
      {\originalsection[#2]{#3}}%
    }%
  \thispagestyle{section}%
}

% Texto del pie de página.
\fancyhead[LE,RO]{\rule[-4ex]{0pt}{2ex}\sffamily\textsl{\rightmark}}
\cfoot{}                                                 % Centro.
\fancyhead[LO,RE]{\sffamily{\leftmark}}
\fancyfoot[LE,RO]{\Large\sffamily\thepage}


% ---------------------------------------------------------------------------
% ENTORNOS PARA MATEMÁTICAS
% ---------------------------------------------------------------------------

% Nuevo estilo para definiciones.
\newtheoremstyle{definition-style} % Nombre del estilo.
{}               % Espacio por encima.
{}               % Espacio por debajo.
{}                   % Fuente del cuerpo.
{}                   % Identación.
{\bf\sffamily}                % Fuente para la cabecera.
{.}                  % Puntuación tras la cabecera.
{.5em}               % Espacio tras la cabecera.
{\thmname{#1}\thmnumber{ #2}\thmnote{ (#3)}}     % Especificación de la cabecera (actual: nombre en negrita).

% Nuevo estilo para notas.
\newtheoremstyle{remark-style} 
{10pt}                
{10pt}                
{}                   
{}                   
{\itshape \sffamily}          
{.}                  
{.5em}               
{}                  

% Nuevo estilo para teoremas y proposiciones.
\newtheoremstyle{theorem-style}
{}                
{}                
{}           
{}                  
{\bfseries \sffamily}             
{.}                
{.5em}           
{\thmname{#1}\thmnumber{ #2}\thmnote{ (#3)}}

% Nuevo estilo para la demostración
%\renewenvironment{proof}{{\itshape \sffamily Demostración \\}}{\hspace*{\fill}\qed}

\makeatletter
\renewenvironment{proof}[1][\proofname] {\par\pushQED{\qed}\normalfont\topsep6\p@\@plus6\p@\relax\trivlist\item[\hskip\labelsep\itshape\sffamily#1\@addpunct{.}]\ignorespaces}{\popQED\endtrivlist\@endpefalse}
\makeatother

% Configuración general de mdframe, los estilos de los teoremas, etc
\mdfsetup{skipabove=12pt,skipbelow=12pt,innertopmargin=12pt,innerbottommargin=4pt}

% Creamos los 'marcos' de los estilos

\mdfdefinestyle{nth-frame}{
	linewidth=2pt, %
	linecolor= 500, % 
	%linecolor=black,
	topline=false, %
	bottomline=false, %
	rightline=false,%
	leftmargin=0pt, %
	innerleftmargin=1em, % 
	rightmargin=0pt, %
	innerrightmargin=0pt, % % 
	%splittopskip=\topskip, %
	%splitbottomskip=\topskip, %
}% 

\mdfdefinestyle{nprop-frame}{
	linewidth=2pt, %
	linecolor= 300, %
	%linecolor= gray, 
	topline=false, %
	bottomline=false, %
	rightline=false,%
	leftmargin=0pt, %
	innerleftmargin=1em, %
	innerrightmargin=0em, 
	rightmargin=0pt, %
	%splittopskip=\topskip, %
	%splitbottomskip=\topskip, %
}%       

\mdfdefinestyle{ndef-frame}{
	linewidth=2pt, %
	linecolor= 300, % 
	%linecolor= gray!50,
	backgroundcolor= 50,
	%backgroundcolor= gray!5,
	topline=false, %
	bottomline=false, %
	rightline=false,%
	leftmargin=0pt, %
	innerleftmargin=1em, %
	innerrightmargin=1em, 
	rightmargin=0pt, % 
	innertopmargin=1.5em,%
	innerbottommargin=1em, % 
	splittopskip=\topskip, %
	%splitbottomskip=\topskip, %
}%               

% Asignamos los marcos a los estilos

\surroundwithmdframed[style=nth-frame]{nth}
\surroundwithmdframed[style=nprop-frame]{nprop}
\surroundwithmdframed[style=nprop-frame]{ncor}
\surroundwithmdframed[style=ndef-frame]{ndef}

% Nuevo estilo para ejemplos.
\newtheoremstyle{example-style}
{10pt}                
{10pt}                
{}                  
{}                   
{\bf \sffamily}              
{:}                 
{.5em}               
{}                   

% Teoremas, proposiciones y corolarios.
\theoremstyle{theorem-style}
\newtheorem{nth}{Teorema}[section]
\newtheorem{nprop}{Proposición}[section]
\newtheorem{ncor}{Corolario}[section]
\newtheorem{lema}{Lema}[section]
% Definiciones.
\theoremstyle{definition-style}
\newtheorem{ndef}{Definición}[section]

% Notas.
\theoremstyle{remark-style}
\newtheorem*{nota}{Nota}

% Ejemplos.
\theoremstyle{example-style}
\newtheorem{ejemplo}{Ejemplo}[section]

% Listas ordenadas con números romanos (i), (ii), etc.
\newenvironment{nlist}
{\begin{enumerate}
    \renewcommand\labelenumi{(\emph{\roman{enumi})}}}
  {\end{enumerate}}

% División por casos con llave a la derecha.
\newenvironment{rcases}
{\left.\begin{aligned}}
    {\end{aligned}\right\rbrace}

% ---------------------------------------------------------------------------
% PÁGINA DE TÍTULO
% ---------------------------------------------------------------------------

% Título del documento.
\newcommand{\subject}{Geometría III}

% Autor del documento.
\newcommand{\docauthor}{Doble Grado de Informática y Matemáticas}

% Título
\title{
  \normalfont \normalsize
  \textsc{Universidad de Granada} \\ [25pt]    % Texto por encima.
  \horrule{0.5pt} \\[0.4cm] % Línea horizontal fina.
  \huge \sffamily\subject\\ % Título.
  \horrule{2pt} \\[0.5cm] % Línea horizontal gruesa.
}

% Autor.
\author{\Large\sffamily{\docauthor}}

% Fecha.
\date{\vspace{-1.5em} \normalsize \sffamily Curso 2016/17}


% ---------------------------------------------------------------------------
% COMIENZO DEL DOCUMENTO
% ---------------------------------------------------------------------------
\usepackage{amsmath}
\begin{document}

\maketitle  % Título.
\newpage
\tableofcontents    % Índice
\vfill
\begin{center}
  \includegraphics{by-nc-sa}  % Licencia.
\end{center}
\newpage


% ----------------------------
% Introducción.
% ----------------------------
\section*{Introducción}
En esta asignatura estudiaremos la geometría afín. Nos centraremos en el plano, el espacio y las figuras que se forman en él.
%%% TODO: completar introducción.
\newpage

\section{El espacio afín}
\begin{ndef}[Espacio afín]
  Sea $E$ un conjunto no vacío. Diremos que $E$ es un \textit{espacio afín} asociado a un espacio vectorial $V$ si existe una aplicación
  \begin{align*}
    \varphi: E \times E & \rightarrow V \\
    (a,b) & \mapsto \varphi(a,b)
  \end{align*}
  que cumpla las siguientes propiedades:
  \begin{nlist}
  \item Si fijamos un punto $a$, la aplicación $\varphi_a$ es biyectiva, o lo que es lo mismo:
    $$\forall a \in E, \ \ v \in V, \ \ \exists! \ b \in E:\varphi(a,b)=v$$
  \item Se tiene la relación de Charles, es decir:
    $$\forall a, b, c \in E \ \ \  \varphi(a,b) + \varphi(b,c) = \varphi(a,c)$$
  \end{nlist}
\end{ndef}
\subsection{El plano afín}

%%% TODO: Decir explícitamente quién es la aplicación \phi en el espacio afín R^2
De ahora en adelante, en el desarollo de estos apuntes consideraremos $\R^2$ como el plano afín. En este plano podemos representar tanto \textit{puntos} como \textit{vectores} (realmente lo haríamos en el plano vectorial). No vamos a entrar a detallar todas las propiedades de los vectores, que vamos a suponer que el lector conoce. \\

\begin{figure*}[h]
	\sffamily
  \centering
  \begin{tikzpicture}
	  	\draw[->] (-0.2,0) -- (4,0) node[right] {$x$};
  		\draw[->] (0,-0.2) -- (0,4) node[above] {$y$};
  		\node [fill=500, inner sep=1pt, scale=1.5, circle] at (2,3) {};
  		\draw [dashed, 300] (0,3) -- (2,3);
  		\draw [dashed, 300] (2,0) -- (2,3);
  		\draw[shift={(0,3)}] (2pt,0pt) -- (-2pt,0pt) node[left] {$3$};
  		\draw[shift={(2,0)}] (0pt,2pt) -- (0pt,-2pt) node[below] {$2$};
  \end{tikzpicture}
  \caption{El punto (2,3) en el plano}
\end{figure*}

Lo que vamos a tratar de hacer en este tema es \textit{transformar} las propiedades de los vectores en propiedades de puntos e intentar relacionar estos dos conceptos.

Dados dos puntos $P$ y $Q$ del plano, podemos transformarlos en un vector llamando $\vec{u}$ al vector que une $P$ y $Q$. Así, $\vec{u} = \overrightarrow{PQ} = Q-P$. De forma analítica podemos observar también que si tenemos un punto $P$, podemos sumarle un vector $\vec{u}$ y obtenemos otro punto $Q$. Hemos visto entonces que podemos sumar tanto dos     vectores como puntos con vectores.

%%% TODO: ¿Insertar un dibujo con la suma de un vector a un punto? ¿Suma de vectores?

\begin{ndef}[Combinación lineal de vectores]
  Si tenemos una familia $u_1,\ldots , u_n$ de vectores y $\alpha_1,\ldots, \alpha_n \in \R$ una familia de escalares, entonces a la expresión
  \[
    \sum_{i=1}^n \alpha_i u_i
  \]
  la llamamos \textit{combinación lineal de vectores}.
\end{ndef}

\begin{ndef}[Razón de vectores proporcionales]
  Dados $u$ y $v$ dos vectores proporcionales ($u = \lambda v$), llamamos \textit{razón} de esos dos vectores al escalar $\lambda \in \R$.
\end{ndef}

A continuación vamos a estudiar distintas \textit{combinaciones} de puntos del plano afín.

\begin{ndef}[Combinación baricéntrica de puntos]
  Dado $\{P_1,...,P_n\}$ un conjunto de puntos del plano afín y $\alpha_1,\dots,\alpha_n \in \R$ un conjunto de escalares que cumplan que $\sum_i \alpha_i = 1$, llamamos \textit{combinación baricéntrica de puntos} a:
  \[
    \sum_{i=1}^n \alpha_i P_i
  \]
  Los escalares $\alpha_1,\dots,\alpha_n$ se llaman \textit{pesos} de los puntos.
\end{ndef}

\begin{nprop} Una combinación baricéntrica de puntos es un punto. 
\end{nprop}

\begin{proof}

  Sea $P = \sum_i \alpha_i P_i$ una combinación baricéntrica de puntos. Para simplificar la demostración, vamos a comenzar transformando los puntos en vectores.

  Restamos un punto cualquiera de la combinación como por ejemplo $P_1$, $P-P_1 = \sum \alpha_i(P_i-P_1)$, obteniendo así una identidad entre vectores de la forma:
  \[
    P-P_1 = \sum_{i=2}^n \alpha_i \overrightarrow{P_1P_i}
  \]

  Y llegamos a que
  \[
    P = P_1 + \sum_{i=2}^n \alpha_i \overrightarrow{P_1P_i}
  \]
  
  Hemos demostrado que $P$ es el resultado de sumar un vector a un punto (en este caso $P_1$), que es otro punto como queríamos.
\end{proof}

\begin{ndef}[Combinación convexa]
  Si en una combinación baricéntrica los pesos $\alpha_i$ asignados a los puntos son todos positivos, la llamaremos \textit{combinación convexa}. \end{ndef}

\begin{ndef}[Centro de gravedad de un conjunto de puntos]
  Si $\{P_1,\dots,P_n\}$ es un conjunto de puntos del plano afín y $\alpha_1,...,\alpha_n$ los \textit{pesos} asociados a estos puntos, llamamos a
  \[
    \sum_{i=1}^n \alpha_iP_i
  \]
  el \textit{centro de gravedad del conjunto de puntos}.
\end{ndef}

\begin{ndef}[Punto medio]
  Sean $P$ y $Q$ dos puntos en el espacio afín. Entonces definimos el punto medio entre $P$ y $Q$ como:
  \[
    M := \dfrac{1}{2} P + \dfrac{1}{2}Q = \dfrac{1}{2}(P+Q)
  \]
  Alternativamente, si $u = \overrightarrow{PQ}$, entonces:
  \[
    M = \dfrac{1}{2}u + P
  \]

\end{ndef}

\subsection{Rectas afines}

En este apartado vamos a describir las \textit{rectas afines}. De ahora en adelante, si se hace mención a una recta sin más, daremos por hecho que se trata de una recta afín.

\begin{ndef}[Recta vectorial]
  Una recta vectorial es un subespacio vectorial de dimensión $1$, es decir, para generarla es necesario un único vector. Podemos denotar la recta vectorial generada por un vector $v$ como $L(v)$.
\end{ndef}

\begin{ndef}[Recta afín]
  Si $P$ es un punto y $v$ un vector, entonces una recta afín $r$ viene dada por
  \[
    r = P+L(v) = \{P + \lambda v : \lambda \in \R\}
  \]
  Para distinguir las rectas vectoriales de las rectas afines, notaremos a $L(v)$ como $\vec{r}$.
\end{ndef}

\begin{ndef}[Variedad de dirección]
  Si $v$ es un \textit{vector director}, llamaremos \textit{variedad de dirección} a la recta generada por este vector, $L(v) = \vec{r}$.
\end{ndef}

\textbf{Propiedades de las rectas afines:}
\begin{itemize}
\item Si $Q$ es un punto y $r$ una recta tales que $Q\in r \implies r = Q +L(v)$ donde v es el vector director de la recta $r$.
\item Si $P,Q$ son dos puntos, entonces por ellos pasa una única recta $r = P +L(\overrightarrow{PQ})$. \\
  \begin{proof}
    (Unicidad) Si la recta $r$ pasa por $P$ y $Q$ y sea $s$ otra recta que pasa por ellos, entonces por definición:
    \[
      s = P+L(\overrightarrow{PQ}) = r
    \]
  \end{proof}
\end{itemize}

\subsubsection{Posición relativa entre rectas}

Podemos clasificar un par de rectas $r$ y $s$ por su posición en el plano afín. Así, decimos que dos rectas son

\begin{itemize}
\item Paralelas, notado $r\parallel s$, si y solo si $\vec{r} = \vec{s}$
\item No paralelas, notado $r\nparallel s$, si y solo si sus vectores directores son linealmente independientes o lo que es lo mismo, $u$ y $v$ forman una base.
\end{itemize}

\begin{nth}
  Sean $r$ y $s$ dos rectas del plano afín, entonces:
  \begin{nlist}
  \item $r \nparallel s\implies r \cap s = A$, un punto
  \item $r \parallel s \implies r = s$ o bien $r \cap s =  \emptyset$
  \end{nlist}

  \begin{proof}
    \begin{nlist}
      % REVIEW: No acabo de ver que demuestre lo que dice (i)
    \item Si $u,v$ son los vectores directores de $r$ y $s$ y $P$ y $Q$ puntos de $r$ y $s$ respectivamente, entonces con los pares $\{u,v\}$ tenemos que $P+\lambda u = Q + \mu v$ con $\lambda, \mu \in \R$
      y si $r = P +L(u)$ y $s=Q+L(v)$, y $\overrightarrow{PQ} = \lambda u - \mu v $ y, por ello, existen únicos $\lambda, \mu$.
    \item Sean $u,v$ los vectores directores de $r$ y $s$. Supongamos que además tienen al menos un punto $A$ en común. Escribimos estas rectas como $r = A +L(u)$ y $s=A+L(v)$. Como $r \parallel s \implies L(r)=L(s)$ descubrimos que $r = A +L(u)=A+L(v) = s$.
   	%TODO: demostrar que cuando no existe el punto A la intersección es vacía 
    \end{nlist}
  \end{proof}
\end{nth}


\subsection{Triángulos}
Los \textit{triángulos} son polígonos que vienen dados por tres puntos no alineados $A, B, C$ llamados \textit{vértices} y las tres rectas que los unen.

\begin {ndef}
  Dos o más puntos $A_1, A_2, ... , A_n $ estan alineados $\iff$ existe una recta $r \subset \mathbb{R}^2\ tal\ que\ A_i \in r\ \forall i \in 1,...,n$
\end{ndef}


\begin{nprop}
  $A,B,C$ son puntos no alineados  $ \iff \overrightarrow{AB}, \ \overrightarrow{AC}$ son linealmente independientes.
  \begin {proof}

    \boxed{\Rightarrow}
    Si los puntos $A$, $B$, y $C$ no están alineados entonces $\nexists r \subset \mathbb{R}^2\ $ tal que $A$, $B$, y $C$ $\in r$, luego $\nexists \vec{v}\ tal\ que\ A,B,C \in r = A + L(\vec{v})$. Por tanto, $\overrightarrow{AB} \neq \alpha \cdot \overrightarrow{AC}, \alpha \in \mathbb{R}\ $ porque si no, sería cierto lo anterior. Así, $\overrightarrow{AB}$ y $\overrightarrow{AC}$ son linealmente independientes.

    \boxed{\Leftarrow} Si $A$, $B$ y $C$ están en la misma recta, entonces la única recta que une puntos A y B y los puntos A y C es la misma, luego $r = A + L(\overrightarrow{AB})  = A + L(\overrightarrow{AC})$, por lo que $\overrightarrow{AB}$ y $\overrightarrow{AC}$ serían linealmente dependientes. Como sabemos que $\overrightarrow{AB}$ y $\overrightarrow{AC}$ no lo son, por lo tanto A, B y C no están alineados.
  \end {proof}

\end{nprop}

% REVIEW: Posibilidad de hacer dibujo de triángulo con los 3 puntos y las 3 rectas



\begin{nth}
  Si consideramos los triángulos $ABC$ y $A_1B_1C_1$ con dos de los lados homólogos paralelos y las razones de los vectores de esos lados son iguales.
  Si las rectas paralelas son: $A\vee C \parallel A_1C_1$ y $A\vee B \parallel A_1B_1$, por lo que las razones son:
  \[
    \frac{\vec{A_1C_1}}{\vec{AC}} = \frac{\vec{A_1B_1}}{\vec{AB}}
  \]
  Entonces, la pareja restante $C\vee B \parallel C_1B_1$ y $\frac{\vec{B_1C_1}}{\vec{BC}} = \frac{\vec{A_1B_1}}{\vec{AB}} = \frac{\vec{A_1C_1}}{\vec{AC}}$\\
  \begin{proof}
    Si los lados son paralelos, entonces los vectores directores son proporcionales.

    Si $w$ (el lado no mencionado del primer triángulo) es $w=-u+v$ y $w_1$(el lado no mencionado del segundo triángulo) es $w_1 = - \lambda u + \lambda v = \lambda (-u+v) = \lambda w$  por lo que tenemos lo que queríamos.
  \end{proof}
\end{nth}

\begin{ndef}[Triángulos semejantes]
  Dos triángulos son semejantes si sus lados son paralelos dos a dos. También son semejantes $\iff$ los ángulos son iguales. También son semejantes si sus lados son proporcionales.
\end{ndef}

\begin{ncor}[Teorema de Thales]
  Si en un triángulo $ABC$ se traza una recta paralela a cualquiera de sus lados, obtenemos un triángulo semejante al triángulo original.
\end{ncor}

\begin{nota}
  Dos triángulos tienen todos los lados paralelos dos a dos $\iff$ los segmentos de los lados son proporcionales dos a dos. Esto es por el teorema anterior.
\end{nota}

\begin{ndef}[Triángulo medio]
  Si $ABC$ es un triángulo y $A'$,$B'$,$C'$ los puntos medios de sus lados, entonces:
  \begin{nlist}
  \item $A'B'C'$ es un triángulo
  \item Los lados homólogos son paralelos y el factor de proporcionalidad es: $-\dfrac{1}{2}$
  \end{nlist}\vspace{0.2cm}
  \begin{proof}

    \begin{nlist}
    \item Se hace viendo que si ABC no están alineados, estos 3 tampoco
    \item Tenemos que: $A' = B/2 + C/2$, $B' = A/2+C/2$, $C' = A/2 + B/2$. Tenemos por tanto que: $\vec{A'B'} = B'-A' = \dfrac{1}{2}(A-B) = -\dfrac{1}{2} \vec{AB}$ y por tanto los lados homólogos son paralelos.

      Además, el factor es el número obtenido en la conversión, es decir $\dfrac{-1}{2}$.
    \end{nlist}
  \end{proof}
\end{ndef}


%% Centros y ejes de un triangulo
\subsubsection{Centros y ejes de un triángulo}

\begin{ndef}[Medianas de un triángulo]
  Las medianas de un triángulo $ABC$ son las rectas que unen un vértice con el punto medio del lado opuesto
\end{ndef}
\begin{ndef}[Baricentro de un triángulo]
  El baricentro $G$ de un triángulo que viene dado por:
  \[
    G = \dfrac{1}{3} A +  \dfrac{1}{3} B +  \dfrac{1}{3} C
  \]
\end{ndef}

\pgfdeclarelayer{0}
\pgfdeclarelayer{1}
\pgfdeclarelayer{2}
\pgfdeclarelayer{3}
\pgfsetlayers{0,1,main,2,3}

\begin{center}
  \sffamily
  \begin{tikzpicture}[scale=2]
    \begin{pgfonlayer}{3}
      \coordinate [label=left:$A$] (A) at (0,0);
      \coordinate [label=right:$B$] (B) at (2,0);
      \coordinate [label=above:$C$] (C) at (1,2);
    \end{pgfonlayer}{3}
    \begin{pgfonlayer}{1}
      \draw (A) -- (B) -- (C) -- cycle;
    \end{pgfonlayer}{1}
    \begin{pgfonlayer}{3}
      \node [fill=500,inner sep=1pt,label=below:$M_{AB}$, scale=1.5, circle] (MAB) at
      ($ (A)!.5!(B)$) {};
    \end{pgfonlayer}{3}
    \coordinate (C1) at ($ (MAB)!1.2!(C)$);
    \coordinate (C2) at ($ (MAB)!-.3!(C)$);
    \draw [dashed, 300, name path=MC](C1) -- (C2);
    \begin{pgfonlayer}{3}
      \node [fill=500,inner sep=1pt,label=left:$M_{AC}$, scale=1.5, circle] (MAC) at ($
      (A)!.5!(C)$) {};
    \end{pgfonlayer}{3}
    \coordinate (B1) at ($ (MAC)!1.2!(B) $);
    \coordinate (B2) at ($ (MAC)!-.3!(B)$);
    
    \draw [dashed, 300, name path=MB] (B1) -- (B2);
    
    \begin{pgfonlayer}{3}
      \node [fill=orange,inner sep=1pt,label=right:$M_{BC}$, scale=1.5, circle] (MBC) at ($
      (B)!.5!(C)$) {};
    \end{pgfonlayer}{3}  
    \coordinate (A1) at ($ (MBC)!1.2!(A) $);
    \coordinate (A2) at ($ (MBC)!-.3!(A)$);
    
    \draw [dashed, 300] (A1) -- (A2);
    
    \path [name intersections={of=MB and MC, by=G}];
    \begin{pgfonlayer}{3}
      \node [fill=700, inner sep=1pt, label=below:$G$, scale=1.5, circle] at (G) {};
    \end{pgfonlayer}{3}
  \end{tikzpicture}
\end{center}

\begin{nth}
  Si $ABC$ es un triángulo en el plano afín, entonces las medianas son concurrentes. Su punto de corte es el baricentro $(G)$. El baricentro triseca el segmento que une un vértice con el punto medio homólogo.


\end{nth}

%%% TODO: Elegir cual os gusta más please.

%% \begin{proof}
%%   Tomamos el vértice $A$ y el punto medio del lado opuesto. Lo que haremos es hacer una de las medianas y ver que el Baricentro pertenece a esta. Si tomamos el vértice $A$ y el punto medio del lado opuesto $A'$. Consideramos la recta $AvA' = Av\frac{C+B}{2}$. Tomando $A$ como un punto y $\vec{AA'}$ como el vector director, representamos la mediana de la forma:
%%   \[
%%     \{A + \lambda (\frac{C+B}{2} - A)\}
%%   \] 
%%   Y en este caso basta tomar $\lambda = \dfrac{2}{3}$ para obtener que el baricentro $G = \dfrac{A+B+C}{3}$ pertenece a esta recta.

%%   Ahora, haciéndolo de forma análoga con los otros dos vértices, comprobamos que el baricentro pertenece a las 3 medianas, por tanto es donde trisecan.
%% \end{proof}
\begin{proof}
  Sea $A'$ el punto medio del lado opuesto a $A$ en un triángulo $ABC$. Sea $A \vee A'$ una de las medianas. Esta es:
  \[
    \{\lambda A + \mu A': \lambda, \mu \in \R \text{ , } \lambda+\mu = 1\}
  \]
  Sabemos que $G = \dfrac{1}{3} A +  \dfrac{1}{3} B +  \dfrac{1}{3} C$ y queremos ver si está en la mediana. Lo que queremos comprobar es que:
  \[
    G = \dfrac{1}{3} A +  \dfrac{1}{3} B +  \dfrac{1}{3} C = \lambda A + \mu A'  
  \]
  
  Es claro que $A' = \left(\frac{1}{2}B + \frac{1}{2}C\right)$. Tomamos $\lambda = \frac{1}{3}$, por lo que $\mu = \frac{2}{3}$ ya que el punto que buscamos se encuentra en la mediana. Sustituyendo estos valores en la ecuación anterior:
  
  $$\dfrac{1}{3}A + \dfrac{2}{3}\left(\dfrac{1}{2}B+\dfrac{1}{2}C\right) = \frac{1}{3}A + \frac{1}{3}B + \frac{1}{3}C = G$$
% TODO: Falta demostrar que el baricentro triseca el segmento que une un vértice con el punto medio homólogo.
\end{proof}

Dado un triángulo $ABC$ y su triángulo medio $A'B'C'$, las medianas de cada lado cortan al punto opuesto y a su homólogo del triángulo medio. Las tres medianas concurren en el baricentro, por lo que este triseca a las medianas. Considerando el punto $A$ del triángulo tenemos que $G = \dfrac{2}{3}A' + \dfrac{1}{3}A$

\begin{ndef}[Mediatriz de un segmento]
  Si $\overline{AB}$ es un segmento, entonces llamamos $m_{AB}$  al lugar geométrico $\mathcal{L}=\{P \in \R^2 : AP = BP\}$ (los puntos del plano que distan lo mismo que A que de B).
\end{ndef}

\begin{center}
  \sffamily
  \begin{tikzpicture}
    \coordinate [label=left:$A$] (A) at (0,0);
    \coordinate [label=right:$B$] (B) at (2,0.5);
    \draw (A) -- (B);
    \coordinate (M) at ($ (A)!.5!(B)$);
    \coordinate (top) at ($ (M) !1! 90:(B)$);
    \coordinate (bottom) at ($ (M) !1! 270:(B)$);
    \draw [dashed, 700] (top) -- (bottom) node[right] {Mediatriz del segmento $\overline{AB}$};
  \end{tikzpicture}
\end{center}

\begin{nota}
  A partir de ahora, usaremos $\mu(A,B)$ para denotar la distancia euclídea entre los puntos $A$ y $B$.
\end{nota}

\textbf{Propiedades:}
\begin{nlist}
\item El punto medio $M \in \mathcal{L}$. 
\item La mediatriz es una recta que viene dada por $\langle\vec{AB},P\rangle  = \dfrac{1}{2}(|B|^2-|A|^2)$ y que es perpendicular al segmento $\overline{AB}$ y pasa por $M$ el punto medio del segmento.\\
  \begin{proof}
    La mediatriz está dada por $L = \{P\in\mathbb{R}^2, \mu(P,A) = \mu(P,B)\iff \mu(P,A)^2 = \mu(P,B)^2\}$. Sea $P(x,y),A(x_1,y_1),B(x_2,y_2)$, entonces $\mu(P,A)^2 = |P|^2+|A|^2-2\langle A,P\rangle  = |P|^2+|B|^2-2\langle B,P\rangle = \mu(P,B)^2 \implies$ $$ \frac{|A|^2-|B|^2}{2} =\langle A-B,P\rangle  = \langle\vec{BA},P\rangle $$
    Donde podemos ver que $\frac{|A|^2-|B|^2}{2}$ es una constante y $\langle\vec{BA},P\rangle $ genera la recta:
    \[
      r: \vec{BA}_1x+ \vec{BA}_2y - \frac{|A|^2-|B|^2}{2} = 0
    \]
    que es perpendicular a $\vec{BA}$.
  \end{proof}
\end{nlist}
\begin{nota}
  $|\vec{AP}|^2 = \mu(A,P)^2\implies \mu(A,P)^2 = 	\left(\sqrt{\langle \vec{AP},\vec{AP}\rangle }\right)^2 =  \langle A-P, A-P \rangle  =^{(1)} \langle A,A-P\rangle -\langle P,A-P\rangle  =  \langle A,A\rangle  - \langle A,P\rangle +\langle P,P\rangle - \langle P,A\rangle   = |A|^2+|P|^2-2 \langle P,A\rangle $. Y en (1) hemos aplicado la bilinealidad del producto escalar.
\end{nota}
\begin{ndef}[Rectas perpendiculares.]
  Dos rectas son perpendiculares si sus vectores directores son ortogonales.
\end{ndef}

\begin{ndef}[Mediatrices de un triángulo]
  Si $ABC$ es un triángulo, cada lado de este triángulo tendrá su mediatriz. Por tanto, tendremos 3 mediatrices para cada triángulo.
\end{ndef}

\begin{nth}
  Las tres mediatrices de un triángulo son concurrentes. Además, concurren en el punto $O$ llamado el circuncentro, el centro de la circunferencia que pasa por $A$,$B$ y $C$.
\end{nth}
\begin{proof}
  Lo primero es ver que dos mediatrices nunca son paralelas. Tomamos $O = m_{AB} \cap m_{AC}$. Por la definición de mediatriz tenemos que $\vec{AO} = \vec{BO}$ y que $\vec{AO} = \vec{CO}$. Es obvio entonces que $\vec{BO} = \vec{CO}$ luego $O \in m_{BC}$.
  % TODO: Completar esta demostración.
\end{proof}


\begin{ndef}[Alturas de un triángulo]
  Son las rectas que pasan por un vértice y son perpendiculares al lado opuesto del triángulo. Estas alturas se cortan en el ortocentro $H$.
\end{ndef}

\begin{nota}
  Las mediatrices de un triángulo $ABC$ son las alturas de $A'B'C'$, por lo que las alturas de $A'B'C'$ son concurrentes.
\end{nota}

\begin{nprop}
  Todo triángulo $ABC$ es el triángulo medio de otro triángulo $A''B''C''$
\end{nprop}
\begin{proof}
  Para ello, tenemos que considerar $A'' = C+ \vec{AB}$. Ídem con $B''$ y $C''$.
\end{proof}
\begin{ndef}
  Todo triángulo $ABC$ tiene un triángulo doble asociado, el cual es el hallado en la proposición.
  % FIXME: ¿Qué es un triángulo doble? ¿Qué proposición?
\end{ndef}


\begin{nprop}
  Si $ABC$ es un triángulo y $A'B'C'$ su triángulo medio, entonces las alturas de $A'B'C'$ son las mediatrices de $ABC$ y $O=H'$
\end{nprop}
\begin{proof}
  Si tomamos $m_{BC}$ podems comprobar que pasa por $A'$ y es perpendicular a $\vec{BC}$, por lo que es perpendicular a $\vec{B'C'}$, por lo que será la altura del triángulo medio que pasa por $A'$. Análogamente ocurre con los $B'$ y $C'$, por tanto, $O=H'$.
\end{proof}


\begin{nprop}Si 2 de los centros del triángulo ABC coinciden, entonces el triángulo es equilátero.
\end{nprop}

%TODO: la dejo a medias, comentada para que no ensucie.
%\begin{proof}
%	Si es equilatero $\implies$ coinciden todos los centros, por lo que veremos que si no es equilatero no coinciden.
%\end{proof}

\begin{ndef}[Homotecia]
  Se define homotecia de centro $A$ y razón $\lambda$ como la aplicación $h_{A,\lambda}(P) = \lambda P + (1-\lambda)A = A+\lambda(\vec{AP})$
\end{ndef}

\begin{nprop} Si $h$ es una homotecia de razón $\lambda\ne 0$, entonces 
\begin{enumerate}
\item $h$ multiplica la distancia por $|\lambda|$.
\item $h$ lleva rectas perpendiculares en rectas perpendiculares.
\item $h$ lleva mediatrices en mediatrices.
\item $h$ lleva alturas en alturas.
\end{enumerate}
\end{nprop}

\begin{proof}\hfill
\begin{enumerate}
	\item $h(P)-h(Q)=(A+\lambda(\vec{AP})) -(A+\lambda(\vec{AQ})) = \lambda(A-P) - \lambda(A-Q) = \lambda (Q-P)$.
	\item  Sean $r$ y $ s$, rectas perpendiculares, $P,P' \in r$ y $Q,Q'\in s$ puntos de las rectas. $\langle h(P) - h(P'),h(Q) - h(Q')\rangle  =  \langle \lambda (\vec{PP'}, \lambda(\vec {QQ'}))=^1 \lambda^2 <\vec{PP'},\vec{QQ'}\rangle  = \lambda^2 \cdot 0 = 0 $, donde en 1 se ha usado la bilinealidad del producto escalar.
	\item Si $h$ multiplica distancia por $\lambda$, las distancias iguales siguen siendo iguales, por lo que el conjunto de puntos que estan a la misma distancia de dos rectas se mantiene.
	\item Como mantiene distancias iguales, mantiene puntos medios. Además al mantener perpendicularidad, la recta perpendicular que pasa por el punto medio se mantiene.
\end{enumerate}
\end{proof}


\begin{nprop}
	
Para toda homotecia, el centro, un punto y su imagen están alineados.
\end{nprop}

\begin{proof}
	Sea la homotecia $h_{A,\lambda}(P) = \lambda P + (1-\lambda)A = A+\lambda(\vec{AP})$ y el punto cualquiera$P$. Veamos que en la recta (única) que contiene a $P$ y a $ A$, $r= A + \mu \vec{AP}$,  contiene a $h_{A,\lambda} (P)$. Esto es obvio por la misma definición de homotecia.
	
	
\end{proof}


\begin{nth}[Recta de Euler]
  Dado un triángulo $ABC$ no equilátero, entonces sus tres centros $G,O,H$ están alineados(están en la recta de Euler).
\end{nth}
\begin{proof}
  Tomamos el triángulo y su triángulo medio $A'B'C'$. Entonces, hay una homotecia de centro $G$ y razón $\dfrac{-1}{2}$ que lleva $ABC$ en $A'B'C'$.(Probarlo, aplicando la fórmula de la homotecia con su cent y su razón). 
  Ahora, si $G',O',H'$ son los centros de $A'B'C'$ y $G,O,H$ los centros de $ABC$, tenemos que $G'=G$ y $O=H'$. Sabemos que la homotecia $h$ lleva $G \to G$, $O \to O'$ y $H \to H' = O$.
  
  Entonces, si aplicamos $h$ al centro G, a un punto $H$ y a su imagen $O$, por el ejercicio anterior, estan alineados.
  
\end{proof}



\subsection{Ecuaciones de una recta}
Hemos visto una recta como un punto $A$ y un vector $\vec{v}$. Es decir, $r \equiv (x_0,y_0) +(x,y) : (x,y) \in L(v) \equiv (x_0+x,y_0+y) $.

También por tanto podemos hacer una recta mediante dos puntos, tomando el vector que forman esos dos puntos y uno de ellos y teniendo la situación anterior.

Las rectas vectoriales en el plano vectorial son siempre de la forma:
\[
  \vec{r} \equiv (x,y) \in L : \alpha x + \beta y = 0 \text{  con $\alpha$ o $\beta$ distinto de 0}
\]

Si $r$ es una recta afín, entonces:
\[
  r \equiv \alpha x + \beta y = \gamma
\] con $\alpha$ o $\beta$ distintos de $0$ y $\alpha,\beta,\gamma \in \R$

Ahora, si tenemos $r'\equiv \alpha' x + \beta' x = \gamma '$ podemos ver que:
\begin{itemize}
\item $r \parallel r' \iff (\alpha,\beta), (\alpha',\beta') $ son linealmente dependientes. Es decir, si el determinante que forman es $0$

\item $r \nparallel r'\iff (\alpha,\beta), (\alpha',\beta') $ no son linealmente dependientes. Es decir, si el determinante que forman no es $0$
\end{itemize}

\subsection{Cuadriláteros}

\begin{ndef}
  $ABCD$ es un cuadrilátero que está formado por 4 puntos consecutivos no alineados 3 a 3. Se llaman lados a los segmentos que unen los vertices consecutivos, y diagonales a los segmentos que unen $A \vee C$ y $B \vee D$.
\end{ndef}

\begin{ndef}[Lado]
  Un lado es una recta que une dos vértices consecutivos. En concreto, los cuadriláteros tienen 4 lados: $A \vee B$, $B \vee C$, $C \vee D$
  y $D \vee A$.
\end{ndef}

\begin{ndef}[Diagonal]
  Una diagonal es una recta que une dos vertices no consecutivos. En concreto, los cuadriláteros tienen 2 diagonales: $A \vee C$ y $B \vee D$.
\end{ndef}

\begin{ndef}[Paralelogramo]
  Un paralelogramo es un cuadrilátero que cumple que todos sus lados opuestos son paralelos. Es decir, sea $ABCD$ un
  paralelogramo entonces $AvB$ es paralelo a $C \vee D$ y $B \vee C$ es paralelo a $D \vee A$.
  

\end{ndef}
\hfill
%TODO : Incluir dibujo de un paralelogramo con sus respectivos
\begin{nprop}
	Los paralelogramos cumplen que
  siendo $V_{1} = A \vee B$, $V_{2} = B \vee C$, $V_{3} = C\vee D$, $V_{4} = D \vee A$ entonces:
  \begin{nlist}
  \item $V_{1} - V_{2} - V_{3} + V_{4} = 0$
  \item $V_{3} = \lambda V_{1}\ con\ \lambda \neq 0$
  \item $V_{4} = \mu V_{2}\ con\ \mu \neq 0$
  \end{nlist}
\end{nprop}

\begin{proof}
	\begin{nlist} 
		\item Como $V_1 = V_3$ y $V_2 = V_4$ al restarlos da 0.
		\item Al ser esos lados paralelos, sus vectores han de ser proporcionales
		\item Análoga al anterior
	\end{nlist}
\end{proof}


\begin{nprop}
  Sea ABCD un cuadrilátero. Entonces, son equivalentes:
  \begin{nlist}
  \item ABCD es un paralelogramo
  \item $\vec{AB} = \vec{DC}$
  \item Las diagonales del cuadrilátero se cortan en su punto medio.
  \end{nlist}
\end{nprop}


\begin{proof}\hfill
  \boxed{1) \Rightarrow 2)}
  Solo tenemos que fijarnos en lo que implica ser paralelogramo (ver propiedad arriba):
  $V_{1} - V_{2} - V_{3} + V_{4} = 0 \Leftrightarrow V_{1} - V_{2} = V_{3} - V_{4} $
  Sustituimos $V_{3} = \lambda V_{1}$ y $V_{4} = \mu V_{2}$:
  $V_{1} - V_{2} = \lambda V_{1} - \mu V_{2}$
  Como sabemos que $V_{1}\ y\ V_{2}$ unen 3 puntos que no están alineados, entonces sabemos que son linealmente independientes y por
  tanto forman base, por esta razon, cualquier vector se forma como combinación lineal de ellos, con coeficientes únicos.
  Concluyendo por tanto que entonces $\lambda = 1 y \mu = 1$

  \boxed{2) \Rightarrow 3)}
  Tenemos que $B - A = C - D$, pasando los términos:
  $B + D = C + A$
  Y ahora dividiendo todo entre 2:
  $\frac{1}{2}B + \frac{1}{2}D = \frac{1}{2}C + \frac{1}{2}A$
  Con lo cual obtenemos los puntos medios de las diagonales que son iguales.

  \boxed{3) \Rightarrow 2)} Del mismo modo, si se cortan en la diagonal 
  $\implies \frac{1}{2}B + \frac{1}{2}D = \frac{1}{2}C + \frac{1}{2}A \implies 
  \frac{1}{2}B - \frac{1}{2}A = \frac{1}{2}C - \frac{1}{2}D \implies $ $B - A = C - D \implies \vec{AB} = \vec{DC}$
  
  \boxed{2) \Rightarrow 1)} Para que sea un paralelogramo, además de $\vec{AB} = \vec{DC}$, ha de cumplir $\vec{AD} = \vec{BC}$, pero por hipotesis $A - B = D - C \implies A - D = B - C$
  
\end{proof}

\begin{ndef}
  \begin{itemize}
  \item Un cuadrilátero es convexo si las dos diagonales separan a los otros vértices.\\
    \begin{tikzpicture}
      \draw (0,0) node {}
      -- (1,3) node {}
      -- (4,2) node {}
      -- (3,0) node {}
      -- cycle;
      \draw [dashed, gray] (0,0) -- (4,2);
      \draw [dashed, gray] (1,3) -- (3,0);
    \end{tikzpicture}
  \item Un cuadrilátero es no convexo si las dos diagonales no separan a los otros vértices.\\
    \begin{tikzpicture}
      \draw (0,0) node {}
      -- (0,4) node {}
      -- (4,0) node {}
      -- (4,3) node {}
      -- cycle;
    \end{tikzpicture}
  \end{itemize}
\end{ndef}

\subsection{Transformaciones}

\begin{ndef}[Afinidades]
  Se define una afinidad f como una aplicacion $f:\R^2 \rightarrow \R^2$ con $f(x) = Ax + b$, con $x =$ $\begin{pmatrix}
    x\\
    y\\
  \end{pmatrix}$, $A =$ $\begin{pmatrix}
    a & b\\
    c & d\\
  \end{pmatrix}$
  y $b =$ $\begin{pmatrix}
    b1\\
    b2\\
  \end{pmatrix}$


\end{ndef}

\begin{ndef}[Traslaciones]
  Se define la traslacion de vector $v$ como $t_v: \R^2 \rightarrow \R^2$ con $t_v(P) = P + v$
\end{ndef}


\begin{nprop}
  Se dan las siguientes propiedades:
  \begin{nlist}
  \item $t_v \circ t_w = t_{v+w}$
  \item $t_{v}^{-1} = t_{-v}$
  \item $t_0 = I$
  \end{nlist}
\end{nprop}
\begin{proof}\hfill \begin{center}\end{center}
  \begin{nlist}
  \item $t_v \circ_w = t_v(P+w) = P+w+v = t_{v+w}(P)$
  \item Vamos a ver que $t_{-v}\circ t_v = I = t_v \circ t_{-v}$
    \begin{itemize}
    \item $t_{-v}\circ t_v (P) = P-v+v = P = I(P)$
    \item $t_v\circ t_{-v} = P+v-v = P = I(P)$
    \end{itemize}
  \item $t_0(p) = P +0 = P = I(P)$
  \end{nlist}
\end{proof}


\begin{nota}
  Durante esta sección, consideraremos continuamente la afinidad $f(P) = AP + b $
\end{nota}


\begin{nprop}
  Se dan las siguientes propiedades:
  \begin{nlist}
  \item Las traslaciones son un caso particular de afinidades, en las que $A = I$ y $b = v$
  \item Sean f y g afinidades, entonces $(g \circ f) (P)$ es afinidad
  \item $f^{-1}(P) = A^{-1}P +A^{-1}b$
    
  \end{nlist}
\end{nprop}
\begin{proof}
  \begin{nlist}
  \item Trivial
  \item $(f\circ g)(P) = f(AP +b) = A'(AP)+A'b +b' $, donde $AA'$ es una matriz y $A'b+b'$ es un vector, por lo que tenemos una nueva afinidad.
  \item Es directo por la composición de funciones.
  \end{nlist}
\end{proof}

\begin{ndef}[Función vectorial asociada]
  Se llama función vectorial asociada a una afinidad($\vec{f}$) a la función que lleva vectores en vectores tal que $\vec{f}(v) = Av$
\end{ndef}

\textbf{Propiedades de las afinidades:}\\
\begin{enumerate}
\item $f(P + \vec{v}) = f(P) + \vec{f}(v)$
\item $f$ lleva rectas en rectas.
\item $f(r\cap s) = f(r) \cap f(s)$
\item $\lambda,\mu \in \R$, $\lambda+\mu=1$. Entonces, $f(\lambda P + \mu Q) = \lambda f(P) +  \mu f(Q)$
\item $f$ lleva puntos medios en puntos medios
\item $f$ lleva segmentos en segmentos
\item $f$ lleva rectas paralelas en rectas paralelas
\item $f$ lleva triángulos en triángulos
\item $f$ lleva medianas en medianas
\item $f$ lleva baricentros en baricentros
\item $f$ lleva cuadriláteros en cuadriláteros
\item $f$ lleva paralelogramos en paralelogramos
\end{enumerate}
\begin{proof}\hfill
  \begin{enumerate}
  \item $f(P+\vec{v}) = A(P+\vec{v}) + b = AP + A\vec{v} + b = AP + b + A\vec{v} = f(P) + \vec{f}(\vec{v})$
  \item Sea $r = P + \lambda \vec{v}$ una recta. Entonces, $f(r) = f(P + \lambda \vec{v}) =^{(1)} f(P) + \lambda \vec{f}(\vec{v})$, donde en $(1)$ hemos usado la linealidad de la asociada vectorial.
  \item Equivalente a la siguiente propiedad al ser $\lambda P + \mu Q$ con $\mu+\lambda = 1$ las combinaciones baricéntricas de $P$ y $Q$ y por lo tanto una recta.
  \item $f(\lambda P + \mu Q) = A(\lambda P + \mu Q) + b =A(\lambda P + \mu Q) + (\lambda + \mu)b = \lambda (AP + b) + \mu (AQ + b) = \lambda f(P) + \mu f(Q)  $
  \item Si $PQ$ es un segmento y $f(\dfrac{1}{2}P + \dfrac{1}{2}Q) = \dfrac{1}{2}f(P) +  \dfrac{1}{2}f(Q)$
  \item Trivial
  \item Trivial
  \item Trivial por las 4 anteriores.
  \item Trivial pues $f$ lleva rectas en rectas.
  \item Trivial pues lleva intersecciones de rectas en intersecciones de rectas
  \item Trivial pues lleva rectas en rectas e intersecciones en intersecciones
  \item Trivial por la anterior y porque mantiene paralelismo

    
  \end{enumerate}
\end{proof}

% \begin{nota}
%   Tenemos que ver esta transformación como una transformación intuitiva en la cual
% \end{nota}



\subsection{Circunferencias}

\begin{ndef}[Circunferencia]
  Una circunferencia es el lugar geométrico de los puntos del plano $\R^2$ que distan exactamente una distancia $r$ de un punto fijo. Las notaremos como $C = C(centro,radio)$
\end{ndef}


%FIXME: ¿es esto cierto? ¿O hay una que leva la primera en la segunda y otra la segunda en la primera?
\begin{nprop}
  Dadas dos circunferencias, existen exactamente dos dilataciones que llevan una en la otra.
  \begin{enumerate}
  \item $C_1 = C(A_1,p_1)$
  \item $C_2 = C(A_2,p_2)$
  \end{enumerate}
  Si $h$ es una dilatación de razón $\lambda$, entonces $h(p) = \lambda p + b$.
  En nuestro caso, cumplirán que $h(C_1) = C_2$ y $p_2 = \lambda p_1$
\end{nprop}


\begin{ndef}
	Una dilatación es una homotecia con $\lambda \ne 1$.
\end{ndef}


\begin{ndef}
  Dos puntos $A$, $B$ de $C = C(O,p) $ se dicen diametralmente opuestos si y solamente si $O \in \bar{AB} \equiv O = \dfrac{1}{2}A + \dfrac{1}{2}B$
\end{ndef}


%%% DONE
\begin{nprop}
  Un par de puntos diametralmente opuestos de $C$ determinan la circunferencia de forma única.
\end{nprop}
\begin{proof}
	Sabemos la existencia de la circunferencia de manera trivial, porque el punto medio $O$ esta a la misma distancia de ambas. Entoces $C_1 = C(O, |\vec{AB}|/2)$ pasa por ambos. \\
	Supongamos que existe otra circunferencia  distinta con el mismo centro que pasa por ambos. Entonces como pasa por uno de los puntos también tiene el mismo radio y por lo tanto son distintas.\\
	Por último debemos ver si el centro puede ser diferente de $O$. Como son diametralmente opuestos el centro debe estar en la linea que los une. Además como deben estar a la misma distancia tiene que ser el punto medio, es decir, $O$.
	
\end{proof}


\begin{ndef}[Ángulo recto]
  Dos vectores forman un ángulo recto si $<\vec{PA},\vec{PB}> = 0$, siendo $P$ el punto donde se cortan.
\end{ndef}

%TODO: completar
\begin{nprop}
  El lugar geométrico desde el que dos puntos dados del plano se ven bajo un ángulo recto es una circunferencia.
\end{nprop}

\begin{proof}
	Estos son todos los puntos $P$ que cumple $<P-A, P-B>=0$. Entonces por la bilienealidad y la comutatividad del producto escalar $<P-A,P-B> = <P-A,P> - < P-A,B> = <P,P> - <P,A> - <P,B> + <A,B> = | P|^2 - <P,A+B> + <A,B> = 0  $. Definimos $M=A/2+B/2 $, entonces, %aquí usamos a aros y su magia. Demostrar. 
	$|P|^2 - <P,2M> + |M|^2 = |P|^2 - 2<P,M> + |M|^2 = |P-M|^2 = \rho$ ,  por lo que la distancia de todos estos puntos a $M$ es igual, por lo que forman una circunferencia de centro $M$ y radio $\rho$. 
\end{proof}

\begin{nota}
	Vease que esta circunferencia es la única definida por $A$ y $B$ de modo que sean diametralmente opuestos.
\end{nota}

%%%TODO:Subir foto de la imagen de la circunferencia de los 9  puntos para poder demostrarlo
\begin{nth}[La circunferencia de los 9 puntos de un triángulo]
  Sea $ABC$ un triángulo en el plano $\R^2$. Entonces existe una circunferencia que pasa por los siguientes puntos:
  \begin{itemize}
  \item $A',B',C'$ los puntos medios de los lados
  \item $A_2,B_2,C_2$ los puntos medios de los segmentos $\bar{HA}$, $\bar{HB}$, $\bar{HC}$
  \item $A_3,B_3,C_3$ los tres pies de las alturas.
  \end{itemize}
\end{nth}
\begin{proof}
  Tenemos que $B'A'B_2A_2$ es un rectángulo
  %%% TODO:COMPLETAR DEMOSTRACIÓN.
\end{proof}



\subsection{Ángulos}

\begin{ndef}[Ángulo entre vectores]
  Dos vectores $\vec{w}$ y $\vec{v}$ en $\R^2$ forman un ángulo:
  \[
    cos \theta =  \frac{< v,w >}{|v||w|} \implies \theta = cos^{-1} \frac{< v,w >}{|v||w|}
  \]
\end{ndef}
\begin{nota}
  Un ángulo entre rectas sería el ángulo que forman sus vectores directores
\end{nota}
\begin{nprop}
  Si tenemos cuatro rectas $r \parallel r'$ y $s \parallel s'$, entonces los ángulos entre $r$ y $s$, y entre $r'$ y $s'$ son iguales.
\end{nprop}

\begin{nprop}
  Los ángulos opuestos por el vértice formados por dos rectas son iguales.
\end{nprop}


\begin{nprop}
  Si $r$ y $s$ son perpendiculares, entonces $\theta = \pi/2$
\end{nprop}

\begin{nth}
  Los tres ángulos de un triángulo suman $\pi$
\end{nth}
\begin{proof}
  La prueba se hace realizando el triángulo medio y trasladando los ángulos a uno de los vértices de este por rectas paralelas.
\end{proof}

\begin{nprop}
  Si $f:\R^2 \to \R^2$ es una dilatación:
  \[
    f \begin{pmatrix}
      x\\
      y 
    \end{pmatrix} = \lambda\begin{pmatrix}
      x\\
      y 
    \end{pmatrix} + b \quad \lambda \ne 0
  \]
  Las distancias se multiplican por $|\lambda|$, los ángulos se mantienen.
\end{nprop}




\subsection{Rectángulos}
\begin{ndef}
  Un rectángulo es un paralelogramo tal que los lados contiguos son perpendiculares.
\end{ndef}


\begin{nprop}
  Un paralelogramo es rectángulo $\iff$ las dos diagonales miden lo mismo.
\end{nprop}
\begin{proof}\hfill \\
  $\boxed{\Rightarrow}$\\
  %%%FIXME:¿HARÍA FALTA EXPLICARLA MÁS? no me salio muy clara.
  Al ser ambas diagonales hipotenusas de tiágulos idénticos (por ser $\vec{AB}=\vec{CD}$, $\vec{BC}=\vec{DA}$ y por ser recto el ángulo que une los segemtos) son de igual distancia.\\
  %TODO:Incluir imagen de un rectángulo ABCD, con la diagonales dibujadas
  $\boxed{\Leftarrow}$ \\
  Por ser los lados opuestos paralelos de la misma distancia, y la diagonal también, ambos triángulos ABC y BCD son iguales. Entonces, además sabemos que  $\angle A = \angle
 B = \angle C  = \angle D$. Por último, como un cuadrilatero está compuesto por dos triángulos, sabemos que sus ángulos deben sumar $2\pi$, por lo que cada ángulo debe
  medir $\pi/2$, por lo que todos los lados son perpeniculares.
  
\end{proof}

\begin{nota}
	Recordemos para la siguiente demostración que el punto de corte de las diagonales de un cuadrilatero es el punto medio de cada una de ellas.
\end{nota}

\begin{nprop}
  Todo rectángulo se inscribe en una circunferencia. Además,cada par de vértices opuestos son diametralmente opuestos.
\end{nprop}

\begin{proof}
	Que se inscribe en una circunferencia es obvio porque porque al medir las dos diagonales lo mismo, el punto donde se cortan está a la misma distancia de los 4 puntos. Además los lados estarán mas cerca de este centro que los vertices. 
	
	Además como el centro está en la diagonal $A\vee C$, estos dos puntos son diametralmente opuestos. Análogo $B\vee D$
\end{proof}

%%% TODO:completar
\begin{nprop}
  Un paralelogramo se puede inscribir en una circunferencia $\iff$ es un rectángulo
\end{nprop}





%TODO:Si alguien puede que envie al grupo apuntes de este día.
\subsection {Razón aurea}
\begin{ndef}[La razón áurea]
  
\end{ndef}











\section{Movimientos}

\begin{ndef}[Movimientos]
  Un movimiento es una aplicación $f:\R^n \to \R^n$ con $n= 2$ ó $3$ que es una afinidad que conserva distancias.
  \[
    |f(P)-f(Q)|= |P-Q| \quad  \quad |\vec{f}(\vec{QP})|= |\vec{QP}|
  \]
\end{ndef}

\begin{nprop}
  $f$ es un movimiento $\iff$ $\mathcal A$ es ortogonal ($\mathcal A ^{-1}= \mathcal A ^t$). Es de la forma:
  \[
    f \begin{pmatrix}
      x_1\\
      \vdots\\
      x_n 
    \end{pmatrix} = \mathcal A \begin{pmatrix}
      x_1\\
      \vdots\\
      x_n 
    \end{pmatrix} + \begin{pmatrix}
      b_1\\
      \vdots\\
      b_n 
    \end{pmatrix}
  \]
\end{nprop}


\begin{ejemplo}
  \begin{nlist}
  \item Traslaciones $p \mapsto p+v$
  \item Reflexión central $p \mapsto -p +b= -p +2A$ de centro $A$
  \item Giro de ángulo $\theta$ y centro $A$ (en $\R^2$).
    \[
      f= \begin{cases}
	f(A) = A \\
	\vec{f} \equiv A_{\theta} = \begin{pmatrix}
          cos \theta & -sen \theta\\
          sen \theta & cos \theta
        \end{pmatrix} 
      \end{cases}
    \]
  \end{nlist}
\end{ejemplo}

\begin{nprop}
  Si realizamos un giro $\mathcal A_\theta$ y sobre ese giro aplicamos otro giro $\mathcal A_{\theta'}$, tenemos que:
  \[
    \mathcal A_{\theta} \circ \mathcal A_{\theta'} = \mathcal A_{\theta+\theta'}
  \]
  Siendo $\mathcal A$ la matriz del giro.
\end{nprop}

\begin{ndef}[Matrices ortogonales]
  Son de la forma:
  \[
    \begin{pmatrix}
      a & -b \\
      b & a 
    \end{pmatrix} \quad \text { ó } \quad \begin{pmatrix}
      a & b \\
      b & -a 
    \end{pmatrix}
  \]
  con $a^2+b^2 = 1$
\end{ndef}

\begin{ndef}[Movimiento directo]
  Es un movimiento cuya matriz es ortogonal. 
\end{ndef}
\begin{nprop}
  En $\R^n$, si $\lambda = 1$ no es un valor propio de $\mathcal A$, entonces $f$ tiene un punto fijo y solo 1.
\end{nprop}

\subsection{Tipos de movimientos}
\begin{itemize}
\item Movimiento deslizante si no hay puntos fijos $f = t_v \circ t_r$. En estas, el vector deslizamiento es $v=b^1$
\item Auténtica reflexión: reflexión con puntos fijos $f=r_r$. En este caso, tenemos que $r= \dfrac{1}{2}b^{-1}+V_1$
\end{itemize}

\subsection{Giros en el espacio}
\begin{ndef}[Movimiento helicoidal]
  Es un movimiento en el que no hay puntos fijos. Su matriz en cierta base ortonormal es:
  \[
    \begin{pmatrix}
      cos \theta & -sen\theta & 0 \\
      sen \theta & cos\theta & 0\\
      0 & 0 & 1
    \end{pmatrix} 
  \]
  En ellas, podemos obtener el ángulo sabiendo que 
  \[
    tr(A) = 1+2cos\theta
  \]
\end{ndef}

%% VOY A COPIAR LO DE CLASSE, REVISAR SI VA AQUÍ

%%% METER AQUÍ APUNTES DE CAMINO

\subsection{Movimientos en $\R^3$}

\[
  f \begin{pmatrix}
    x\\
    y\\
    z
  \end{pmatrix} = A\begin{pmatrix}
    x\\
    y\\
    z
  \end{pmatrix} + \begin{pmatrix}
    b_1\\
    b_2\\
    b_3
  \end{pmatrix}
\]

\begin{itemize}
\item Si la matriz $A$ es simétrica, $f$ es una reflexión que puede ser respecto de un plano o respecto de una recta (axial o especular). Además, tenemos el caso de que sea deslizante o no deslizante.

  \begin{itemize}
  \item Simetría axial, $dim(V_1) = 1$, $dim(V_{-1}) =2$. Su determinante vale por tanto 1. Es un movimiento directo.
  \item Simetría especular , $dim(V_1) = 2$, $dim(V_{-1}) =1$. Su determinante vale -1.Es un movimiento inverso.
  \end{itemize}

\item Si la matriz $A$ no es simétrica, entonces podemos tener un movimiento directo con $dim(V_1) = 1$ y $dim(V_{-1})=0$, teniendo por tanto un giro o un movimiento helicoidal.
  
\item Movimientos inversos, en los que hay puntos fijos por fuerza y tenemos que $dim(V_1) = 0$ y $dim(V_{-1}) = 1$. Este movimiento lo llamaremos giro con simetría. La matriz de este movimiento, en ciertas coordenadas, es:
  \[
    \begin{pmatrix}
      cos \theta & -sen\theta  & 0\\
      sen\theta & cos\theta  & 0\\
      0 & 0 & -1
    \end{pmatrix} 
  \]
\end{itemize}

\begin{ejemplo}
  Obtener la expresión matricial en coordenadas usuales de la simetría respcto del plano $x-y+z = 1\equiv \Pi$
\end{ejemplo}
\begin{proof}[Solución:]
  Tomamos un punto $P(x,y,z)$ y tomamos su simétrico respecto del plano $\sigma(P)$. Si tomamos el vector normal al plano, el $n=(1,-1,1)$, sacado de la ecuación del plano, podemos determinar una recta normal al plano de la forma:
  \[
    s\equiv Q = P + \lambda n = (x,y,z)+\lambda(1,-1,1)
  \]
  Ahora, tenemos que buscar la intersección de la recta con el planno, introducimos la ecuación de la recta en el plano. Si desarrollamos la ecuación nos queda $(x+\lambda,y-\lambda,z+\lambda)$ y si lo introducimos en el plano:
  \[
    x-y+z+3\lambda = 1 \implies \lambda = \frac{1-x+y+z}{3}
  \]
  Por tanto, el punto de corte del plano con la recta estará en el punto:
  \[
    Q = (x,y,z)+\frac{1-x+y+z}{3}(1,-1,1) \in s\cap \Pi
  \]
  Ahora, tratamos de hallar $\sigma(P)$. Para ello, como el punto de intersección de la recta con el plano sería el punto medio del vector $\vec{PQ}$ , podemos hacer :
  \[
    \sigma(P) = P + 2 \vec{PQ} = (x,y,z) + \dfrac{2}{3}(1-x+y+z)(1,-1,1)
  \]
  Y ya tendríamos, para cada punto $P$, su simétrico. Ahora, desarrollamos la expresión y la transformamos en expresión matricial.
\end{proof}

\begin{nth}
  Dados $ABC, A'B'C'$ dos triángulos $\Rightarrow \exists f$ una afinidad que lleva uno a otro. Son equivalentes:
  \begin{enumerate}
  \item Los triángulos son semejantes.
  \item Los lados homólogos son proporcionales.
  \item Los ángulos homólogos son iguales.
  \end{enumerate}	
\end{nth}

\begin{proof} \hfill 

  $\boxed{1 \Rightarrow 2}$ Trivial por definición de semejanza.\\
  $\boxed{1 \Rightarrow 3}$ Trivial, ya que como la semejanza es la composición de homotecia y movimiento y tanto las homotecias como los movimientos conservan los ángulos.\\
  $\boxed{2 \Rightarrow 1}$ Se $h$ una homotecia de razón $\lambda = \frac{a'}{a}$. Llamemos $A''B''C'' = h(ABC)$. Entonces, $A'B'C'$ y $A''B''C''$ tienen sus lados homólogos de la misma logitud, luego existe $g$ movimiento que lleva cada vértice de $A'B'C'$ en el correspondiente de $A''B''C''$ Entonces $g^{-1} \circ h$ lleva $ABC$ en $	A'B'C'$. Como $g^{-1}$ es un movimiento y $h$ una homotecia, $g^{-1} \circ h$ es una semejanza.
\end{proof}


\begin{nprop}[Propiedad del ángulo inscrito]
  Si $C$ es una circunferencia de centro $O$ y $A$ un punto cualquiera de ella y B y C dos puntos distintos y distintos de $A$, entonces el ángulo que forman los segmentos $AB$,$AC$ es justo la mitad del que forman $OB$, $OC$. Esto es:
  \[
    A = \dfrac{1}{2} O
  \]
\end{nprop}
\begin{ncor}
  Todos los ángulos inscritos que abarcan el mismo arco son iguales.
\end{ncor}

\begin{ndef}[Arco Capaz]
	Es el lugar geométrico de los puntos del plano $\R^2$ desde los que un segmento dado $\bar{AB}$ se observa bajo un ángulo dado.
\end{ndef}

\begin{ejemplo}
	El arco capaz de un segmento $\bar{AB}$ es la unión de dos carcos de circunferencia con extremos $A$ y $B$ simétricos respecto de $AvB$.
	
	\begin{proof}[Solución]
	Por la propiedad el ángulo inscrito, los dos arcos están contenidos en $\mathcal L$. Sea 
	%%%% COMPLETARRRRRRRRRRR
	
\end{proof}
\end{ejemplo}
\begin{ejemplo}
	Sea $ABCD$ un cuadtrilátero convexo. $ABCD$ está inscrito en una circunferencia $\iff $\^{A}$ +$ \^{C}$ = \pi$ y \^{B} + \^{D} $ = \pi$
	\begin{proof}
	Supongamos que no, supongamos que $ABC \in \mathcal C$. Ahora, tomamos $D' \in \mathcal C$ tal que $ABCD' \in \mathcal C $
\end{proof}
\end{ejemplo}

\begin{ndef}[Triángulo órtico]
	SI $ABC$ es un triángulo con con ángulos menores que $\pi/2$,entonces los pies de las alturas pertenencen al interior de los lados. El triángulo órtico es entonces el triángulo $DEF$ en el que los lados están formados por los puntos donde se cortan las alturas con los lados del inicial.
\end{ndef}

\textbf{Propiedades:}\\
\begin{nlist}
	\item Existe una circunferencia por $ADHF$. Lo mismo para $BEHD$ y $CFHE$
\end{nlist}

\section{Cónicas y Cuádricas}
Ya sabemos que la ecuación de una recta viene dada por una ecuación de la forma:
\[
ax + by +c = 0
\]

Ahora, la ecuación de una elipse vendría dada por una ecuación de una forma:
\[
ax^2 + by^2 + c = 0
\]
También la ecuación de una hipérbola viene dada por una ecuación de la fora;
\[
axby = c
\]
También si multiplicamos dos ecuaciones de una recta, tendríamos la ecuación de una cónica en $\R^2$
\[
(ax + by +c)(a'x+b'y+c') = 0
\]

\subsection{Cambio de coordenadas en el espacio afín.}

En el espacio afín $\R^3$ tomaremos dos sistemas de referencia ortonomales:
\begin{nlist}
	\item El usual: con $\mathcal R = \{O, e_1, e_2, e_3\}$ con $O=(0,0,0)$, $e_i$ tiene un 1 en la posición $i$.
	\item El formado por otro origen $\bar O = (a,b,c)$ y una base ortonormal $\{\bar e_1, \bar e_2, \bar e_3\}$. 
\end{nlist}

La matriz de cambio de base es una matriz ortogonal $P$ cuyas columnas son las coordenadas de $\bar e_i$ con $i=1,2,3$ respecto de la base usual.

Así, el cambio de coordenadas para vectores será:
\[
\begin{pmatrix}
x\\
y\\
z
\end{pmatrix}  = P \begin{pmatrix}
\bar x\\
\bar y\\
\bar z
\end{pmatrix} 
\]

Y para puntos será:
\[
\begin{pmatrix}
x\\
y\\
z
\end{pmatrix} = P \begin{pmatrix}
\bar x\\
\bar y\\
\bar z
\end{pmatrix}  + \begin{pmatrix}
 a\\ 
 b\\
 c
\end{pmatrix} 
\]





\subsection{Cu\'adricas con centro en $\mathbb{R}^3$.}


Una cuadr\'ica  $\cal Q$ es el lugar geom\'etrico de los puntos $P=(x,y,z)\in\mathbb{R}^3$ tales que
\[ 
\alpha_1 \,x^2 + 2\alpha_2 \,x y + 2\alpha_3\,x z + \alpha_4 \,y^2 + 2 \alpha_5 \,y z + \alpha_6 \,z^2 
+ 2\beta_1 \,x+ 2\beta_2 \,y +2 \beta_3 \,z + \gamma = 0
\]
o en forma matricial 
\[  
\left( \begin{array}{ccc} 
x & y & z   
 \end{array} \right)\begin{pmatrix}
 \alpha_1 & \alpha_2 / 2 &\alpha_3 / 2\\
 \alpha_2 / 2 & \alpha_4 &\alpha_5 / 2\\
 \alpha_3 / 2& \alpha_5 / 2& \alpha_6
\end{pmatrix} 
\left(\begin{array}{c}
x \\
y \\
z
\end{array} \right)
+
 2\left( \begin{array}{ccc}
\beta_1 & \beta_2 & \beta_3   
 \end{array} \right)
\left( \begin{array}{c}
x \\
y \\
z
 \end{array} \right)
+ \gamma = 0
\]

donde $\cal A$ es una matriz $3\times 3$ sim\'etrica y no nula y $\beta_i, \gamma\in\mathbb{R}$. 

\begin{nota}
	Recordemos que la traza de una matriz es igual a la suma de los elementos de su diagonal.
\end{nota}

\begin{nprop}
	Toda matriz simétrica es diagonalizable por una matriz ortogonal:
	\[
	\mathcal P ^t \mathcal A \mathcal P = \bar{ \mathcal A}, \quad \text{$\bar{\mathcal A}$ es simétrica}
	\]
\end{nprop}
\begin{nota}
	Recordemos que diagonalizar por semejanza es hacer $P^{-1}AP$ u hacerlo por conjugación es hacer $P^t A P$
\end{nota}

\begin{nth}
	Si una $A$ es una matriz simétrica, entonces $\exists = \{\bar e_1, \bar e_2, \bar e_3\}$ una base ortonormal de vectores propios
\end{nth}

\begin{nprop}
	El centro de una cuádrica $\mathcal Q$ es un punto $P=(a,b,c)$ tal que la reflexión central $\sigma_A: P \mapsto -P+2A$ se tiene que $\sigma_A (\mathcal Q) = \mathcal Q$
\end{nprop}


\begin{ndef}[Centro]
	Un {\it centro} de $\cal Q$ es un punto $A=(a,b,c)$ que es un centro de simetr\'ia de $\cal Q$, es decir, tal que la reflexi\'on central $\sigma_A:\mathbb{R}^3\longrightarrow\mathbb{R}^3$, $\sigma_A(P)= -P+2A$ lleva  $\cal Q$ en $\cal Q$. Hacemos el cambio de coordenadas que toma a $A$ como nuevo origen y 
y una base ortonormal tal que la matriz del cambio de base sea $\cal P$.
\end{ndef}


Expresamos la cu\'adrica en las nuevas coordenadas
\[ 
\left( \begin{array}{ccc} 
\overline{x} & \overline{y} & \overline{z}   
 \end{array} \right)\overline{{\cal A}}
\left(\begin{array}{c}
\overline{x} \\
\overline{y} \\
\overline{z}  
\end{array} \right)
+
2
 \left( \begin{array}{ccc}
\overline{\beta_1} & \overline{\beta_2} & \overline{\beta_3}   
 \end{array} \right)
\left( \begin{array}{c}
\overline{x} \\
\overline{y} \\
\overline{z}  
 \end{array} \right)
+ \overline{\gamma} = 0
\]
siendo
\[
\overline{{\cal A}} = {\cal P}^t{\cal A}{\cal P} \hspace{1.5cm}
\left( \begin{array}{ccc}
\overline{\beta}_1 & \overline{\beta}_2 & \overline{\beta}_3   
 \end{array} \right) = 
  \left\{
\left( \begin{array}{ccc} 
a &  b &  c 
 \end{array} \right){\cal A} +
 \left( \begin{array}{ccc}
\beta_1 & \beta_2 & \beta_3   
 \end{array} \right) \right\}
{\cal P}
\]


\[
 \overline{\gamma} = \left( \begin{array}{ccc} 
a &  b & c 
 \end{array} \right){\cal A}
\left(\begin{array}{c}
a  \\
b \\
c  
\end{array} \right)  + 
2 \left( \begin{array}{ccc}
\beta_1 & \beta_2 & \beta_3   
 \end{array} \right)
\left( \begin{array}{c}
a \\
b \\
c  
 \end{array} \right)
+ \gamma
\]

\vspace{.3cm}

Para que el origen de las nuevas coordenadas sea un centro de la c\'uadrica imponemos que la parte lineal de la ecuaci\'on anterior se anule

\

\[
{
\left( \begin{array}{ccc} 
a &  b &  c 
 \end{array} \right){\cal A}
+
 \left( \begin{array}{ccc}
\beta_1 & \beta_2 & \beta_3   
 \end{array} \right)
=0 }
\]
y la ecuaci\'on de la cu\'adrica en esas coordenadas queda
\[
{\cal Q} \, \, \, \, \equiv \, \, \, \,
\left(\begin{array}{ccc} 
\overline{x} & \overline{y} & \overline{z}   
 \end{array} \right){\overline{\cal A}}
\left(\begin{array}{c}
\overline{x} \\
\overline{y} \\
\overline{z}  
\end{array} \right) + \overline{\gamma} = 0
\]
siendo $\overline{\gamma}$ la constante dada por
\[
\overline{\gamma} = 
\left\{
\left( \begin{array}{ccc} 
a &  b &  c 
 \end{array} \right){\cal A}
+
 2\left( \begin{array}{ccc}
\beta_1 & \beta_2 & \beta_3   
 \end{array} \right)\right\}
\left( \begin{array}{c}
a \\
b \\
c  
 \end{array} \right)
+\gamma = 
\left( \begin{array}{ccc}
\beta_1 & \beta_2 & \beta_3   
 \end{array} \right)
\left( \begin{array}{c}
a \\
b \\
c  
 \end{array} \right)+\gamma
\]
\[
{ 
\overline{\gamma} =
\beta_1 a+ \beta_2 b+ \beta_3  c 
+\gamma
}
\]

\subsubsection{\bf Ecuaci\'on reducida de una cu\'adrica con centro}

Para obtener la ecuaci\'on reducida de la cu\'adrica, consideramos el sistema de referencia que tiene como origen el centro de $\cal Q$ y como vectores una base ortonormal de vectores propios de la matriz $\cal A$ y valores propios $\lambda_1,\lambda_2,\lambda_3$. Entonces
\[
{\cal R} =\{A; {e}_1, {e}_2, {e}_3 \}
\hspace{2cm}
{\cal Q} \, \, \, \equiv \, \, \, \lambda_1 {x}^2 +\lambda_2 {y}^2 +\lambda_3 {z}^2 +
{\gamma} = 0.
\]
A continuaci\'on damos la lista de los distintos tipos de cu\'adricas con centro y  sus ecuaciones reducidas. Cuando ${\gamma}\neq 0$, estas ecuaciones se obtienen dividiendo la ecuaci\'on anterior por ${\gamma}$. 

\begin{itemize}
	
\item
{ {\bf Rango de ${\cal A} = 3$.}} 

\vspace{.3cm}

Si $\cal A$ tiene rango $3$ entonces $\cal Q$ tiene un \'unico centro.
Entonces existen $a,b,c>0$ tales que en esas coordenadas la cu\'adrica se escribe 


\begin{tabular}{ll}
{\bf  Elipsoide} \hspace{5cm} & 
$
\frac{{x}^2}{a^2} + \frac{{y}^2}{b^2}+\frac{{z}^2}{c^2} = 1
$
\vspace{.5cm}
\\
{\it Elipsoide imaginario}  &
$
-\frac{{x}^2}{a^2} - \frac{{y}^2}{b^2} - \frac{{z}^2}{c^2} = 1
$
\vspace{.5cm}
\\
{\bf  Hiperboloide de una hoja}  &
$
\frac{{x}^2}{a^2} + \frac{{y}^2}{b^2}-\frac{{z}^2}{c^2} = 1
$
\vspace{.5cm}
\\
{\bf  Hiperboloide de dos hojas} &
$
\frac{{x}^2}{a^2} - \frac{{y}^2}{b^2}-\frac{{z}^2}{c^2} = 1
$
\vspace{.5cm}
\\
{\bf  Cono}  &
$
\frac{{x}^2}{a^2} + \frac{{y}^2}{b^2}-\frac{{z}^2}{c^2} = 0
$
\vspace{.5cm}
\\
{\it Cono imaginario}  &
$
\frac{{x}^2}{a^2} + \frac{{y}^2}{b^2}+\frac{{z}^2}{c^2}=0
$
\vspace{.4cm}
\end{tabular}



\item {\bf Rango ${\cal A}=2$.} 


Consideramos cu\'adricas con centro y rango $2$ en el sistema de referencia $\cal R$

\vspace{.5cm}

\begin{tabular}{ll}
{\bf Cilindro el\'iptico} \hspace{5cm} & 
$
\frac{{x}^2}{a^2} + \frac{{y}^2}{b^2}=1
$
\vspace{.5cm}
\\
{\it Cilindro imaginario}  &
$
-\frac{{x}^2}{a^2} - \frac{{y}^2}{b^2}=1
$
\vspace{.5cm}
\\
{\bf  Cilindro hiperb\'olico}  &
$
\frac{{x}^2}{a^2} - \frac{{y}^2}{b^2}=1
$
\vspace{.5cm}
\\
{\bf  Par de planos reales}  &
$
\frac{{x}^2}{a^2} - \frac{{y}^2}{b^2}=0
$
\vspace{.5cm}
\\
{\it Par de planos imaginarios} &
$
\frac{{x}^2}{a^2} + \frac{{y}^2}{b^2}=0
$
\vspace{.3cm}
\\
\end{tabular}
 


\item {\bf Rango ${\cal A}=1$.} 

\vspace{.4cm}

Cu\'adricas con centro y rango $1$ en el sistema de referencia $\cal R$

\vspace{.5cm}

\begin{tabular}{ll}
{\bf  Par de planos paralelos.} \hspace{5cm} & 
$
\frac{{x}^2}{a^2}=1
$
\vspace{.4cm}
\\
{\it Par de planos imaginarios paralelos}  &
$
\frac{{x}^2}{a^2} =-1
$
\vspace{.4cm}
\\
{\bf  Plano doble}  &
$
{{x}^2} = 0
$
\vspace{.3cm}
\\
\end{tabular}

\end{itemize}


\subsection{\bf  C\'onicas con centro en $\mathbb{R}^2$.}

\vspace{.5cm}

Dada una c\'onica $\cal C$ con centro consideramos el sistema de referencia 
${\cal R}=\{A; e_1,e_2\}$ donde $A$ 
es el centro y $e_1,e_2$ es una base ortonormal de vectores propios de la matriz $\cal A$. Se tienen los siguientes casos

\vspace{.8cm}

\begin{tabular}{ll}
{\bf Elipse} \hspace{5cm} & 
$
\frac{{x}^2}{a^2} + \frac{{y}^2}{b^2} = 1
$
\vspace{.4cm}
\\
{\it Elipse imaginaria}  &
$
-\frac{{x}^2}{a^2} - \frac{{y}^2}{b^2} = 1
$
\vspace{.4cm}
\\
{\bf  Hip\'erbola}  &
$
\frac{{x}^2}{a^2} - \frac{{y}^2}{b^2} = 1
$
\vspace{.4cm}
\\
{\bf  Par de rectas reales concurrentes}  &
$
\frac{{x}^2}{a^2} - \frac{{y}^2}{b^2}=0
$
\vspace{.4cm}
\\
{\it Par de rectas imaginarias} &
$
\frac{{x}^2}{a^2} + \frac{{y}^2}{b^2}=0
$
\vspace{.4cm}
\\
{\bf  Par de rectas reales paralelas}  &
$
\frac{{x}^2}{a^2}=1
$
\vspace{.4cm}
\\
{\it Par de rectas imaginarias paralelas}  &
$
-\frac{{x}^2}{a^2} = 1
$
\vspace{.4cm}
\\
{\bf  Recta doble}  &
$
{{x}^2} = 0
$
\vspace{.3cm}
\\
\end{tabular}
\subsubsection{\bf  Elementos notables de las c\'onicas y cu\'adricas con centro.}

\begin{ndef}[Centro]
	Es el centro de simetría de la cónica, es decir, el punto $A$ tal que la reflexión $\sigma_A$ mantiene invariante la cónica. Es importante notar que no todas las gáficas tienen centro.
\end{ndef}


\begin{ndef}[Ejes]
	 Son los ejes de simetr\'ia. Cada eje es una recta que pasa por el centro y tiene como vector director a un vector propio de la matriz $\cal A$. En $\mathbb{R}^2$ (respectivamente $\mathbb{R}^3$) se consideran dos (respectivamente tres)  ejes perpendiculares pasando por el centro.
\end{ndef}

\begin{ndef}[Vértices]
	La intersecci\'on de los ejes con la c\'onica o cu\'adrica.
\end{ndef}


\subsection{\bf Cu\'adricas sin centro en $\mathbb{R}^3$.}


Sea $\cal Q$ una cu\'adrica sin centro tal que el rango de $\cal A$ sea $2$. Entonces los valores propios son 
$\lambda_1,\lambda_2 \neq 0$ y $\lambda_3=0$. Tomamos el sistema de referencia ${\cal R} = \{O; e_1,e_2,e_3\}$ con una base de vectores propios y la ecuaci\'on queda de la forma 

\vspace{.3cm}

\[
{\cal Q} \, \, \equiv \, \, \lambda_1 x^2 + \lambda_2 y^2 + \beta_1 x + \beta_2 y + \beta_3 z + \gamma = 0
\]  

\vspace{.2cm}

Consideramos un nuevo sistema de referencia $\overline{{\cal R}} = \{A;e_1,e_2,e_3\}$ de manera que el origen 
$A=(\alpha_1,\alpha_2,\alpha_3)$ nos permita obtener cuadrados perfectos en $\overline{x}=x-\alpha_1$ e 
$\overline{y}=y-\alpha_2$. La nueva ecuaci\'on es
\[
{\cal Q} \, \, \equiv \, \, \lambda_1 \overline{x}^2 + \lambda_2 \overline{y}^2 + 
\overline{\beta}_3 \overline{z} = 0
\]
donde $\overline{\beta}_3\neq 0$ ya que la cu\'adrica no tiene centro. Esto nos permite obtener las ecuaciones reducidas para el caso sin centro.


\subsubsection{\bf Ecuaciones reducidas para cu\'adricas sin centro.}

\begin{itemize}

\item {Rango $\cal A$ $=2$.}



{\bf Paraboloide El\'iptico.}    \hspace{2.7cm}   $z = \frac{x^2}{a^2} + \frac{y^2}{b^2}$

\vspace{.2cm}

{\bf Paraboloide Hiperb\'olico.}    \hspace{2cm}   $z = \frac{x^2}{a^2} - \frac{y^2}{b^2}$


\begin{tikzpicture}
\sffamily
\begin{axis}
\addplot3 [surf,mesh/interior colormap=
           {inside}{color=(50) color=(300)},colormap/blackwhite,shader=flat,draw=black] {x^2-y^2};
\end{axis}
\end{tikzpicture}

\vspace{0.5cm}

\item {Rango $\cal A$ $=1$.}


{\bf Cilindro Parab\'olico}    \hspace{3.3cm}   $z = \frac{x^2}{a^2}$ 

\end{itemize}

\subsection{\bf C\'onicas sin centro en $\mathbb{R}^2$.}



{\bf Par\'abola} \hspace{5.5cm}  $y = \frac{x^2}{a^2}$



\subsection{Cuádricas sin centro}

\section{El plano proyectivo}

Para definir correctamente el plano proyectivo, primero realizaremos una "definición" a través de un caso concreto. Situémonos en $\mathbb{R}^3$. Ahora definamos el plano $\pi \equiv z = 1 \equiv$ plano afín, asi como el plano $\pi_0 \equiv z = 0 \equiv$ plano vectorial, y el punto $O = (0,0,0)$ el origen.


Entonces definimos el plano proyectivo $\mathbb{P}^2$ como el conjunto formado por las intersecciones de cada recta vectorial de $\mathbb{R}^3$ con el plano $\pi$. Este conjunto contiene a todos los puntos del plano $\pi$(dado un punto del plano siempre puedes encontrar una recta centrada en el origen que pase por el), pero ademas, contiene la interseccion de las rectas vectoriales paralelas a dicho plano. Por ello llamaremos "puntos" del plano proyectivo a cada intersección de una recta vectorial con $\pi$ y en concreto "puntos en el infinito"\ aquellos puntos generados por una recta paralela a $\pi$. Simbólicamente:
$$\mathbb{P}^2 = \mathbb{R}^2\ \cup \{puntos\ del\ infinito\}$$
Ahora definimos la rectas de $\mathbb{P}^2$ como la intersección de un plano vectorial con $\pi$, de manera que de la misma que forma que ocurria con los puntos, $\mathbb{P}^2$ contiene a todas las rectas de $\mathbb{R}^2$ y a un recta más, que es la definida por el plano $\pi_0 \equiv z= 0$ a la que llamaremos recta del infinito ($r_\infty$). Si nos fijamos, esta recta contiene a todos los puntos del infinito, por tanto:
$$\mathbb{P}^2 = \mathbb{R}^2\ \cup r_\infty$$

\begin{nth}
Sean P y Q dos puntos distintos de $\mathbb{P}^2$, existe una única recta proyectiva que contiene a P y Q. A esta la llamaremos $P \vee Q$
\end{nth}
\begin{proof}

De lo estudiado en espacio afín sabemos que dos rectas vectoriales distintas determinan un plano vectorial unico que las contiene, por tanto, utilizando la biyeccion que conocemos entre rectas vectoriales y puntos de $\mathbb{P}^2$ y planos vectoriales y rectas de $\mathbb{P}^2$, la prueba del resultado es trivial.
\end{proof}

\begin{nth}
Dos rectas proyectivas se cortan en único  punto proyectivo
\end{nth}
\begin{proof}
La demostracion es similar a la del teorema anterior usando que dos planos vectoriales distintos se cortan un una única recta vectorial
\end{proof}

\textbf{Consequencias:}

Una recta $r$ en $\mathbb{R}^2$ identifica a una única recta en $\mathbb{P}^2$ que contiene a:
\begin{nlist}

	\item $r \cap \mathbb{R}^2 = la\ recta\ de\ siempre$ 
%%%TODO Pensar una forma de escribirlo mejor
	\item $r \cap {puntos\ en\ el\ infinito} = r \cap r_\infty = 1\ punto\ en\ el\ infinito$

\end{nlist}

Ahora, una vez comprendida la aproximación inicial necesaria para poder imaginarnos y comprender el plano proyectivo, veamos que se puede construir con dos planos cualesquiera con una serie de condiciones.


Sean dos planos $\pi'$ y $\pi_0'$, cumpliendo que:
\begin{nlist}
	\item $O \notin \pi'$
	\item $O \in \pi_0$
	\item $\pi\ ||\ \pi_0$ 
\end{nlist}
Con estos dos planos podemos realizar la misma construcción que realizamos al comienzo del tema y formando de nuevo el plano proyectivo $\mathbb{P}^2$. 

Si nos fijamos, podemos decir que el plano proyectivo formado por cada par de planos es el mismo porque los elementos de ambos estan determinados inequívocamente por las rectas vectoriales y los planos vectoriales, por tanto, ambos son iguales, con la diferencia de que la llamada "recta en el infinito" es distinta, que estaria determinada respectivamente por los planos $\pi_0$ y $\pi_0'$

\begin{nprop}
Sea $r \in \mathbb{P}^2 \implies \mathbb{P}^2 - r = plano\ afin$
\end{nprop}
%%%TODO Esto hace falta demostrarlo? En dicho caso, como se haría?

\begin{ndef}
Tres puntos no alineados $A,B,C \in \mathbb{P}^2$ definen unico triángulo $ABC$, cuyos elementos son similares a los definidos para los triángulos en el plano afín
\end{ndef}

\begin{ndef}
Cuatro puntos no alinedos 3 a 3 $A,B,C,D \in \mathbb{P}^2$ definen un único cuadrilatero, para el que definiremos los siguientes elementos:
\begin{itemize}
	\item Vértices: $A, B, C, D$
	\item Lados: $A\vee B, B\vee C, C\vee D, D\vee A$
	\item Puntos diagonales: $P = (B\vee C) \cap (D\vee A)$, $Q = (A\vee B) \cap (C\vee D)$ 
	\item Rectas diagonales: $A\vee C, B\vee D$
\end{itemize} 
\end{ndef}

\begin{nth}[Teorema de Fano]
	Dada un cuadrilátero $ABCD \in \mathbb{P}^2$, los dos puntos diagonales $P$ y $Q$ y el punto de corte R de las dos rectas diagonales $R = (A \vee C)\cap (B\vee D)$, no están alineados 
\\
	

\end{nth}
\begin{proof}
	Para llevar a cabo esta demostracion trabajaremos en el plano afín $\mathbb{P}^2 -r$, con $r = P\vee Q$. De esta manera obtenemos que el cuadrilátero $ABCD$ que obtenemos en dicho plano afín es un paralelogramo (pues sus lados son paralelos por estar sus puntos de corte en el infinito). 
	
	Por tanto, como sabemos del estudio del plano afín que las diagonales de un paralelogramos se cortan en su punto medio, $R$. Consecuentemente, $P$ y $Q$ pertenecen a $r$ y R al plano afín $\mathbb{P}^2 - r \implies R \notin r$ luego $P,Q$ y $R$ no están alineados.  
	\end{proof}

\begin{nprop}
	Si $r$, $s$ son dos rectas en $\mathbb P^2$ con $r\ne s$, entonces:
	\[
	r \cap s = 1 \text { punto } \begin{cases}
	\in \R^2 \equiv \text{ no son paralelas}\\

	\in r_\infty \equiv \text{ son paralelas}
	
\end{cases}
	\]
\end{nprop}

\begin{nth}[Teorema de Desargues]
	Sean $ABC$, $A'B'C'$ dos triángulos en $\mathbb P^2$ cuyos vértices homólogos son diferentes y los lados homólogos también tales que los lados homólogos se cortan en tres puntos alineados.\\
	Entonces, las rectas que unen los vértices homólogos son concurrentes
\end{nth}

Tenemos también versiones afines del teorema de Desargues:

\begin{nth}[Versión afín del Teorema de Desargues]
	Sean dos triángulos en el afín tales que los lados homólogos no son paralelos y se cortan en tres puntos alineados.\\
	  Entonces las rectas que unen los vértices homólogos son o bien concurrentes o bien paralelas
\end{nth}

\begin{nth}[Versión equivalente 2]
	Sean $ABC$,$A'B'C'$ dos triángulos en el afín tales que los pares de lados homólogos son paralelos.\\
	Entonces, las rectas que unen los vértices homólogos son concurrentes o son paralelas.
\end{nth}


%%%TODO: Definir Q. 
\begin{nth}[Versión equivalente 3]
	En el afín $\mathbb P^2 - P\vee Q$ con $P$ el punto de corte de las rectas que unen los vértices homólogos y $Q$  el punto de corte, si un paralelogramo está inscrito en un cuadrilátero.\\
	Entonces, la otra diagonal del cudrilátero es paralela al otro par de lados del paralelogramo.
\end{nth}

\begin{nota}
	Los enunciados afines son equivalentes al enunciado proyectivo. Por tanto, probando una de las versiones habremos probado todas.
\end{nota}
\begin{proof}
	Vamos a probar la versión afín 2.\\
	Los lados homólogos de los triángulos $ABC$, $A'B'C'$ de $\R^2$ son paralelos. Entonces, existe una dilatación:
	\[
	f: \R^2 \to \R^2 \begin{cases}
	A \mapsto A'\\
	B \mapsto B'\\
	C \mapsto C'
	
\end{cases}
	\]
	Que lleva uno en otro.
	\begin{enumerate}
	\item Si $f$ es una traslación, $f = t_v$, entonces $\vec{AA'}= \vec{BB'} = \vec{CC'}=v$
	\item Si $f$ es una homotecia, $f= h_{p,\lambda}$, con $\lambda \ne 0, 1$, entonces $(P,A,A')$, $(P,B,B')$, $(P,C,C')$ están alineados y $A\vee A'$, $B \vee B'$ y $C \vee C'$ concurren.
\end{enumerate}
\end{proof}

\begin{ndef}[Hexágono]
	Un hexágono son 6 puntos consecutivos y distintos y lados(rectas que unen vértices consecutivos)
\end{ndef}
\begin{nth}[Teorema de Pappus]
	Si los vértices de un hexágono en $\mathbb P^2$ están alternativamente sobre 2 rectas distintas, entonces los 3 pares de lados opuestos se cortan en 3 puntos alineados
\end{nth}

\begin{nth}[Versión afín del teorema de Pappus]
	Dado un hexágono en el plano afín cuyos vértices están alternativamente en dos rectas distintas $s$ y $s'$. Si los pares de lados opuestos son paralelos, entonces el tercer par de lados opuestos también es paralelo.
\end{nth}
%%%% TODO: Incluir dibujos para aclarar las rectas.
\begin{proof}[Demostración de la versión afín]\hfill
\begin{nlist}
	\item Si $s \parallel s'$
\end{nlist}
\end{proof}
\end{document}
